\documentclass[12pt,a4paper]{article}

%====================================================
% BASIC PACKAGES  (same setup as Report 3)
%====================================================
\usepackage[table]{xcolor}
\usepackage{float}
\usepackage{tabularx}
\usepackage{enumitem}


\setcounter{secnumdepth}{5}
\setcounter{tocdepth}{5}

\usepackage{tabularx}
\usepackage{array}
\usepackage{makecell}
\usepackage{longtable}

\newcolumntype{L}[1]{>{\raggedright\arraybackslash}p{#1}}
\newcolumntype{C}[1]{>{\centering\arraybackslash}p{#1}}
\newcolumntype{M}[1]{>{\centering\arraybackslash}m{#1}}

\newcolumntype{Y}{>{\raggedright\arraybackslash}X}

%====================================================
% ENCODING & LANGUAGE
%====================================================
\usepackage[T1]{fontenc}
\usepackage[utf8]{inputenc}
\usepackage[vietnamese,english]{babel}

%====================================================
% FONT
%====================================================
\usepackage{mathptmx}

%====================================================
% MATH
%====================================================
\usepackage{amsmath}
\usepackage{amsfonts}
\usepackage{amssymb}

%====================================================
% GRAPHICS
%====================================================
\usepackage{graphicx}
\usepackage{pdflscape}

%----- diagram placeholder macro ---------------------
% \kfig{width}{path WITHOUT extension}{caption}{label}
% Prefers the vector <path>.pdf exported from the draw.io SVG
% (stays sharp at any zoom level), falls back to the raster
% <path>.png, and finally to a placeholder box so the document
% always compiles.
\newcommand{\kfig}[4]{%
\begin{figure}[H]
    \centering
    \IfFileExists{#2.pdf}{%
        \includegraphics[width=#1,height=0.9\textheight,keepaspectratio]{"#2.pdf"}%
    }{%
        \IfFileExists{#2.png}{%
            \includegraphics[width=#1,height=0.9\textheight,keepaspectratio]{"#2.png"}%
        }{%
            \setlength{\fboxsep}{10pt}%
            \fbox{\parbox[c][5.5cm][c]{0.82\textwidth}{\centering
                \textit{Diagram pending export}\\[6pt]
                \texttt{#2.pdf}}}%
        }%
    }
    \caption{#3}
    \label{#4}
\end{figure}}

%====================================================
% TABLES
%====================================================
\usepackage{booktabs}
\usepackage{longtable}
\usepackage{tabularx}
\usepackage{array}
\usepackage{multirow}
\usepackage{caption}

%====================================================
% PAGE LAYOUT
%====================================================
\usepackage{geometry}

\geometry{
    top=2.5cm,
    bottom=2.5cm,
    left=3.0cm,
    right=2.0cm
}

%====================================================
% SPACING
%====================================================
\usepackage{setspace}
\onehalfspacing

\usepackage{indentfirst}
\setlength{\parindent}{1.27cm}
\setlength{\parskip}{6pt}

%====================================================
% SECTION FORMAT
%====================================================
\usepackage{titlesec}

%====================================================
% PARAGRAPH NUMBERING
%====================================================

\makeatletter

\renewcommand{\theparagraph}
{\thesubsection.\arabic{paragraph}}

\@addtoreset{paragraph}{subsection}

\def\toclevel@paragraph{4}

\renewcommand{\thesubparagraph}
{\thesubsubsection.\arabic{subparagraph}}

\@addtoreset{subparagraph}{subsubsection}

\def\toclevel@subparagraph{4}

\makeatother

\titleformat{\paragraph}
{\normalfont\normalsize\bfseries}
{\theparagraph}
{1em}
{}

\titlespacing*{\paragraph}
{0pt}
{3.25ex plus 1ex minus .2ex}
{1em}

\titleformat{\subparagraph}[runin]
{\normalfont\normalsize\bfseries}
{\thesubparagraph}
{1em}
{}

\titlespacing*{\subparagraph}
{0pt}
{3.25ex plus 1ex minus .2ex}
{1em}

%====================================================
% TABLE OF CONTENTS
%====================================================
\usepackage{tocloft}

\renewcommand{\cftsecleader}{\cftdotfill{\cftdotsep}}
\renewcommand{\cftsubsecleader}{\cftdotfill{\cftdotsep}}
\renewcommand{\cftsubsubsecleader}{\cftdotfill{\cftdotsep}}

\renewcommand{\cftsecfont}{\bfseries}
\renewcommand{\cftsecpagefont}{\bfseries}

\setlength{\cftbeforesecskip}{6pt}
\setlength{\cftbeforesubsecskip}{2pt}

\setlength{\cftparaindent}{\cftsubsubsecindent}
\setlength{\cftparanumwidth}{\cftsubsubsecnumwidth}

%====================================================
% HYPERLINKS
%====================================================
\usepackage{hyperref}

\hypersetup{
    colorlinks=true,
    linkcolor=black,
    citecolor=black,
    filecolor=black,
    urlcolor=blue
}

%====================================================
% COLORS
%====================================================
\definecolor{fptred}{RGB}{192,0,0}
\definecolor{fptBlue}{RGB}{0,0,139}
\definecolor{fptOrange}{RGB}{255,140,0}
\definecolor{headerbg}{RGB}{242,226,204}

%====================================================
% CUSTOM COLUMN TYPES
%====================================================
\newcolumntype{R}[1]{>{\raggedleft\arraybackslash}p{#1}}

%====================================================
% CLASS-SPECIFICATION TABLES
%====================================================
\definecolor{tableHeaderOrange}{RGB}{255,140,0}

% Inline code that tolerates underscores and other TeX special characters.
% Breaks are allowed before camel-case capitals and after punctuation so long
% class names, paths, and error factories do not overflow narrow table cells.
\ExplSyntaxOn
\cs_new_protected:Npn \__code_char:n #1
  {
    \regex_match:nnTF { [A-Z] } {#1}
      { \allowbreak #1 }
      {
        #1
        \regex_match:nnT { [^A-Za-z0-9] } {#1} { \allowbreak }
      }
  }
\NewDocumentCommand{\code}{m}
  {
    \texttt
      {
        \tl_set:Nx \l_tmpa_tl { \tl_to_str:n {#1} }
        \tl_replace_all:Nnn \l_tmpa_tl { ~ } { \c_space_token }
        \tl_map_function:NN \l_tmpa_tl \__code_char:n
      }
  }
\ExplSyntaxOff

% Method table: No | Method | Description
\newenvironment{ClassMethodTable}{%
    \begingroup
    \small
    \setlength{\tabcolsep}{4pt}%
    \renewcommand{\arraystretch}{1.18}%
    \begin{longtable}{|C{0.9cm}|L{5.1cm}|L{9.0cm}|}
    \hline
    \rowcolor{tableHeaderOrange}\textbf{No} & \textbf{Method} & \textbf{Description} \\ \hline
    \endfirsthead
    \hline
    \rowcolor{tableHeaderOrange}\textbf{No} & \textbf{Method} & \textbf{Description} \\ \hline
    \endhead
    \hline
    \endfoot
}{%
    \end{longtable}%
    \endgroup
}

% Attribute table: No | Attribute | Type
\newenvironment{ClassAttributeTable}{%
    \begingroup
    \small
    \setlength{\tabcolsep}{4pt}%
    \renewcommand{\arraystretch}{1.18}%
    \begin{longtable}{|C{0.9cm}|L{8.0cm}|L{6.1cm}|}
    \hline
    \rowcolor{tableHeaderOrange}\textbf{No} & \textbf{Attribute} & \textbf{Type} \\ \hline
    \endfirsthead
    \hline
    \rowcolor{tableHeaderOrange}\textbf{No} & \textbf{Attribute} & \textbf{Type} \\ \hline
    \endhead
    \hline
    \endfoot
}{%
    \end{longtable}%
    \endgroup
}

% Generic three-column supporting table (DTOs, errors, entities).
% The third column absorbs the remaining text width so the table never
% protrudes into the right margin. The third argument is kept for call-site
% readability; its requested width is superseded by the computed width.
\newenvironment{SupportingTable}[3]{%
    \begingroup
    \small
    \setlength{\tabcolsep}{4pt}%
    \renewcommand{\arraystretch}{1.15}%
    \begin{longtable}{|L{#1}|L{#2}|L{\dimexpr\textwidth-6\tabcolsep-1.6pt-#1-#2\relax}|}
    \hline
    \rowcolor{tableHeaderOrange}\textbf{Class / Element} & \textbf{Attributes or Signature} & \textbf{Responsibility} \\ \hline
    \endfirsthead
    \hline
    \rowcolor{tableHeaderOrange}\textbf{Class / Element} & \textbf{Attributes or Signature} & \textbf{Responsibility} \\ \hline
    \endhead
    \hline
    \endfoot
}{%
    \end{longtable}%
    \endgroup
}

%====================================================
% DOCUMENT
%====================================================
\begin{document}

%====================================================
% COVER PAGE
%====================================================
% ========================================================
% FILE: cover.tex
% Cover page of the template (same layout as Report 1 / Report 3):
%   logo, "Capstone Project Report", "Report 4 -- Software Design Document",
%   place and date. The group members and the supervisor are documented in
%   Report 1, exactly as in the other reports of the template.
% ========================================================
\begin{titlepage}
    \centering
    \thispagestyle{empty}
    \hypersetup{pageanchor=false}

    % Logo
    \includegraphics[width=7.5cm]{images/school_logo.png}

    \vspace{2.8cm}

    {\fontsize{28}{34}\selectfont\bfseries\textcolor{fptred}{Capstone Project Report}}\\[0.5cm]
    {\fontsize{22}{27}\selectfont\bfseries\textcolor{fptred}{Report 4 -- Software Design Document}}

    \vspace*{5cm}

    {\fontsize{14}{17}\selectfont -- Can Tho, October 2026 --}

    \vfill

    \hypersetup{pageanchor=true}
\end{titlepage}


%====================================================
% TABLE OF CONTENTS
%====================================================
% ========================================================
% FILE: table_of_contents.tex
% ========================================================

\newpage

\renewcommand{\contentsname}{\fontsize{20}{24}\selectfont\bfseries Table of Contents}

\setcounter{tocdepth}{3} % show sections, subsections and subsubsections

\tableofcontents

\newpage


%====================================================
% MAIN CONTENT
%====================================================

% ========================================================
% FILE: sections/1.Overall_Description.tex
% ========================================================
\section{Overall Description}

\subsection{Assumptions}

This system is designed based on the following assumptions:
\begin{itemize}[leftmargin=*]
    \item The Node.js runtime and the frameworks chosen for the project (NestJS 11 for the backend and Next.js 16 for the frontend) will continue to be actively supported and maintained, ensuring long-term stability and compatibility with evolving technologies.
    \item The RESTful API architecture will remain a scalable and efficient choice for separating the Next.js presentation layer from the NestJS business layer.
    \item PostgreSQL, accessed through Prisma ORM 7, will continue to provide reliable, scalable, and secure relational data storage suitable for handling user identity, jobs, contracts, escrow payments, and forum data.
    \item Third-party services used by the platform --- Stripe (payments and escrow), Cloudinary (file storage), Google (OAuth login), the AI moderation/analysis service, and mail delivery via Nodemailer --- will stay available and API-compatible during the project lifetime.
    \item Modern web browsers (e.g., Chrome, Edge, Firefox, Safari) will continue to support the front-end standards required by Next.js/React 19, including HTML5, CSS3 (TailwindCSS), JavaScript ESM, and responsive rendering with React Server Components.
    \item The application will be deployed in environments with stable internet access to support real-time messaging (Socket.io) and online payment processing.
    \item Users (Guests, Clients, Freelancers, Experts, and Admins) are familiar with basic web application navigation: browsing job boards, filling forms, managing profiles, and following payment flows.
\end{itemize}

This system design is also based on the following risks:
\begin{itemize}[leftmargin=*]
    \item Any breaking change in the major versions of Next.js, NestJS, or Prisma could require refactoring of the user interface components or of the data access layer.
    \item Cross-browser inconsistencies or outdated device software may affect rendering or the real-time messaging experience.
    \item Cybersecurity threats such as credential stuffing, token hijacking, or database breaches could compromise sensitive personal and financial information if not proactively managed (the design therefore includes bcrypt-hashed passwords, JWT access/refresh tokens, session revocation, and role-based access control).
    \item Outages or policy changes of the integrated third-party services (Stripe, Cloudinary, Google OAuth, AI moderation) could temporarily disrupt payments, file uploads, login, or content moderation.
    \item Network connectivity issues could disrupt access to job proposals, contract collaboration, or checkout flows that rely on server interaction.
\end{itemize}

\subsection{Design Constraints}

This system should comply with the following constraints:

\subsubsection{Platform Compatibility}

The Frevia application must be compatible with modern web browsers (e.g., Chrome, Edge, Firefox, Safari) for all user roles including Guest, Client, Freelancer, Expert, and Admin. The frontend is implemented with Next.js 16 (App Router, React Server Components) styled with TailwindCSS 4 and shadcn/ui components; the backend is a NestJS 11 REST API. Responsive web design principles must be applied so that a consistent user experience is provided across desktop, tablet, and mobile browsers. Internationalization is achieved with next-intl so the interface can be served in multiple languages from the same code base.

\subsubsection{Integration Complexity}

The Frevia platform integrates a variety of components: the Next.js web application, the NestJS REST API, the PostgreSQL database (via Prisma 7 with the pg adapter), Stripe Checkout and Stripe Connect for escrow payments and payouts, Cloudinary for image and document storage, Socket.io for real-time chat and notifications, Google OAuth for federated login, e-mail for OTP delivery, and AI services for content moderation, CV parsing, and identity document verification. Communication between the client and the server is carried out through versioned REST endpoints secured with JWT access tokens and refresh tokens stored in dedicated session records; role-based access control restricts each operation to the allowed roles. Design decisions must keep these integrations loosely coupled behind services and repositories so that a provider can be replaced without rewriting business logic.

\subsubsection{Performance Requirements}

The application must remain responsive under concurrent usage by many freelancers and clients. Job browsing, proposal listing, forum feeds, and conversation histories must load promptly; therefore the design mandates server-side pagination for every list operation, database indexes on frequently filtered columns, optimized Prisma queries executed inside transactions for counted lists, and streaming React Server Components to reduce time-to-first-byte. Real-time features (messaging, notifications) must propagate events with minimal delay via Socket.io. Payment operations must reconcile eventual states through Stripe webhooks and explicit synchronization endpoints.

\subsection{Technology Suggestion}

Frevia adopts a RESTful API architecture to support secure, stateless, and scalable communication between the Next.js frontend and the NestJS backend. This approach provides:
\begin{itemize}[leftmargin=*]
    \item Modular, reusable API endpoints for account management, job posting, proposals, contracting, milestone escrow, dispute resolution, community forum, and administration.
    \item Easy extension for mobile applications or third-party integrations in the future.
    \item Secure access using JWT tokens with role-based authorization for every endpoint.
    \item Uniform use of standard HTTP methods (GET, POST, PATCH, PUT, DELETE), making the system intuitive to maintain.
    \item A clear separation between the presentation layer and the business logic layer.
\end{itemize}
All request/response payloads are validated and typed with shared Zod schemas packaged in \texttt{@shared/types}, which is consumed by both the frontend and the backend so that a single source of truth prevents contract drift between the two applications.

\kfig{0.9\linewidth}{images/architecture/restful-api}{RESTful API communication between the Next.js frontend and the NestJS backend}{fig:restful-api}

\subsubsection{NestJS Framework}

The backend of Frevia is built with NestJS 11, a TypeScript-first, modular web framework running on Node.js. It is selected for its:
\begin{itemize}[leftmargin=*]
    \item Rich dependency-injection and module ecosystem, matching the Controller--Service--Repository layering.
    \item First-class decorator-based REST routing with guards and interceptors for authentication, roles, permissions, and response serialization.
    \item Excellent support for Zod schema validation (nestjs-zod) and typed DTOs.
    \item Component structure for WebSockets (Socket.io gateways) used by the messaging module.
    \item Straightforward integration with Prisma ORM and PostgreSQL for persistence.
\end{itemize}

\subsubsection{Next.js and React 19 Frontend}

The web client uses Next.js 16 with the App Router. React Server Components render public pages (job lists, job details, forum) directly from the server, while TanStack Query drives interactive islands (forms, likes, real-time chat). react-hook-form together with Zod resolvers validate user input with the same schemas the API uses.

\subsubsection{Prisma ORM and PostgreSQL}

Persistence relies on PostgreSQL accessed through Prisma 7. The schema defines 62 models covering identity and access, profiles, jobs, proposals, contracts and milestones, payments and disputes, reviews, messaging, the community forum, and platform configuration. Prisma's type-safe client keeps data access code verifiable at compile time.

\subsubsection{Supporting Services}

Payments and payouts use Stripe Checkout and Stripe Connect with webhook-driven reconciliation; binaries are stored in Cloudinary; emails (OTP, notifications) are sent with Nodemailer; Google login uses OAuth 2.0; environment configuration is loaded through Infisical so no secret is committed to the repository.

% ========================================================
% FILE: sections/2.System_Architecture_Design.tex
% ========================================================
\section{System Architecture Design}

\subsection{Overall Architecture}

Frevia follows a three-tier, deployment-separated architecture (Figure~\ref{fig:arch:overall}):

\begin{itemize}[leftmargin=*]
    \item \textbf{Presentation tier --- Next.js 16 web application:} server components render public pages (job board, job detail, community forum, expert directory) and private dashboards; client components handle interactivity such as forms, likes, and chat. All server access goes through a single typed HTTP client that attaches the JWT access token and refreshes it transparently when it expires.
    \item \textbf{Application tier --- NestJS REST API:} a modular monolith organized into 31 route modules. Every module is internally layered as Controller $\rightarrow$ Service $\rightarrow$ Repository. Controllers validate and serialize with shared Zod DTOs; services hold business rules and authorization checks; repositories encapsulate all Prisma queries.
    \item \textbf{Data tier --- PostgreSQL via Prisma 7:} the relational store for 62 domain models. Binary files are never stored in the database but in Cloudinary; real-time transport for chat and notifications is handled by a Socket.io gateway beside the REST API.
\end{itemize}

\kfig{0.95\textwidth}{images/architecture/overall-architecture}{Overall architecture of the Frevia system}{fig:arch:overall}

The two applications are deployed independently (the web app on Vercel, the API on a Node.js host) and communicate only through the documented REST contract validated by the shared Zod schemas package.

\subsection{System Architecture}

The internal architecture of the API follows a strict layered pattern, enforced identically in every module (Figure~\ref{fig:arch:system}):

\begin{itemize}[leftmargin=*]
    \item \textbf{Controller layer} --- thin Routing Layer. Exposes REST endpoints, applies guards (authentication, roles, permissions), binds validated request DTOs, and delegates to a service. Controllers never query the database directly.
    \item \textbf{Service layer} --- business logic layer. Implements use cases, ownership/authorization checks, state-machine transitions (proposal, contract, dispute), and cross-aggregation logic; it throws typed HTTP exceptions (\texttt{NotFoundException 404}, \texttt{ForbiddenException 403}, \texttt{ConflictException 409}) that a global exception filter converts to a uniform error response.
    \item \textbf{Repository layer} --- data access layer. The only component allowed to touch Prisma; complex list queries (pagination + counting) are wrapped in \texttt{\$transaction} for consistency.
    \item \textbf{Shared kernel} --- the \texttt{@shared/types} workspace package contains Zod models, constants, permission names, and error codes shared with the web app.
\end{itemize}

Security flow: every request carries a JWT access token; the auth guard validates it and resolves the active role; role and permission guards authorize module-specific operations; refresh tokens are stored server-side in a \texttt{Session} row and can be revoked from the session management feature.

AI-assisted pipeline: user-generated content (forum posts/comments, CVs, identity documents) is passed through AI moderation or AI analysis services; forum content can be auto-approved when clean, or queued in \texttt{ModerationQueue} for administrator approval when suspicious.

\subsubsection{API Module Structure}

Table~\ref{tab:modules} summarizes the 31 route modules of the API and their responsibilities.

\begin{longtable}{|L{4.2cm}|L{11.0cm}|}
\caption{API route modules and responsibilities}
\label{tab:modules}
\\
\hline
\textbf{Module} & \textbf{Responsibility} \\ \hline
\endfirsthead
\hline
\textbf{Module} & \textbf{Responsibility} \\ \hline
\endhead
auth & registration, OTP, login, refresh token, logout, forgot password, Google OAuth, current account, role switching/joining \\ \hline
sessions & device session listing, detail, revocation \\ \hline
account-profile & general account profile, avatar, identity verification upload/AI check, client profile, social links, favorite freelancers, following \\ \hline
users & administrator user management and skill catalog \\ \hline
roles, permissions & RBAC role and permission administration \\ \hline
profiles & professional profile, CV upload/AI analysis, skills \\ \hline
portfolios & portfolio item CRUD on profiles \\ \hline
expert-profile & expert profile management and public expert directory \\ \hline
profile-revisions & profile change request approval workflow \\ \hline
identity-verifications-admin & identity verification review \\ \hline
job-categories & public job category browsing \\ \hline
job-categories-admin & job category administration (soft delete + restore) \\ \hline
skills-admin & skill catalog administration (soft delete + restore) \\ \hline
browse-job & public job listing and job detail by slug \\ \hline
manage-job & client job management, bookmarks, skill search \\ \hline
proposals & proposal lifecycle: draft, submit, withdraw, interview, reject, accept \\ \hline
contract & contract lifecycle: create, update, sign, complete, cancel \\ \hline
milestone & milestone management inside a contract \\ \hline
milestone-file, milestone-submission & milestone attachments and work submission review \\ \hline
shared-file & contract shared workspace files \\ \hline
payments & Stripe Connect onboarding, platform fee, milestone funding, release, refund, dispute fee, transactions, webhook \\ \hline
disputes & dispute lifecycle (user side and admin arbitration side) \\ \hline
reviews & contract reviews and responses \\ \hline
conversations & direct messaging with real-time delivery \\ \hline
notifications & notification feed management \\ \hline
forums & forum categories and posts (public side) \\ \hline
forums-comment, forums-like, forums-reports & comments, likes, content reports \\ \hline
forums-admin & AI moderation queue, pending posts, trash restore, forum statistics \\ \hline
saved-search & saved job search filters \\ \hline
banners, banners-admin & advertisement banner delivery and administration \\ \hline
\end{longtable}

\kfig{0.95\textwidth}{images/architecture/system-architecture}{Layered system architecture with security and AI moderation flows}{fig:arch:system}

\subsection{Package Diagram}

The repository is a pnpm monorepo managed by Turborepo with three workspaces (Figure~\ref{fig:arch:packages}):

\begin{itemize}[leftmargin=*]
    \item \textbf{\texttt{apps/web}} --- the Next.js application. Sub-packages group screens by actor (guest browsing, client, freelancer, admin), shared UI components (\texttt{@repo/ui}), server actions, and the i18n message catalogs used by next-intl.
    \item \textbf{\texttt{apps/api}} --- the NestJS application. \texttt{src/routes/*} contains the 31 domain modules described above; \texttt{src/shared} contains cross-cutting infrastructure: authentication guards and decorators, permission constants, global HTTP exception filter, Zod serializer interceptor, Cloudinary storage service, and Prisma client setup. Each domain module folder is organized as \texttt{<module>.controller.ts}, \texttt{<module>.service.ts}, \texttt{<module>.repository.ts}, DTO classes, and module-specific error factories (e.g., \texttt{ forums.error.ts}).
    \item \textbf{\texttt{packages/shared-type}} --- the single source of truth for API contracts: Zod request/response models (\texttt{model/*}), enumeration/constant definitions, and permission name lists, imported by both \texttt{apps/web} and \texttt{apps/api}.
\end{itemize}

Dependencies only point downwards: \texttt{web} and \texttt{api} both depend on \texttt{shared-type}; no route module imports another route module --- cross-module data needs flow through services at the Prisma schema level.

\kfig{0.9\textwidth}{images/architecture/package-diagram}{Package diagram of the Frevia monorepo}{fig:arch:packages}

% =================================================================
% FILE: sections/3.System_Detailed_Design.tex
% =================================================================
\section{System Detailed Design}

% =================================================================
% FILE: sections/kvy/sprint1.tex -- Sprint 1 (Community Forum)
% Features are listed in the approved feature order. PNG diagrams are
% rendered next to the .svg exports in images/sprint1/kvy/.
% =================================================================

% ================= View Forum Category List =================
\subsection{View Forum Category List}

\textbf{a.~Class Diagram}

\kfig{0.92\textwidth}{images/sprint1/kvy/classDiagram/View Forum Category List}{Class diagram of View Forum Category List}{fig:cls-view-forum-category-list}

\textbf{b.~Sequence Diagram(s)}

\kfig{0.78\textwidth}{images/sprint1/kvy/sequences/ViewForumCategories}{Sequence diagram of View Forum Category List}{fig:seq-view-forum-category-list}

% ================= View Forum Category Detail =================
\subsection{View Forum Category Detail}

\textbf{a.~Class Diagram}

\kfig{0.92\textwidth}{images/sprint1/kvy/classDiagram/View Forum Category Detail}{Class diagram of View Forum Category Detail}{fig:cls-view-forum-category-detail}

\textbf{b.~Sequence Diagram(s)}

\kfig{0.78\textwidth}{images/sprint1/kvy/sequences/ViewForumCategoryDetail}{Sequence diagram of View Forum Category Detail}{fig:seq-view-forum-category-detail}

% ================= View Post List =================
\subsection{View Post List}

\textbf{a.~Class Diagram}

\kfig{0.92\textwidth}{images/sprint1/kvy/classDiagram/View Post List}{Class diagram of View Post List}{fig:cls-view-post-list}

\textbf{b.~Sequence Diagram(s)}

\kfig{0.78\textwidth}{images/sprint1/kvy/sequences/PostList}{Sequence diagram of View Post List}{fig:seq-view-post-list}

% ================= View Post Detail =================
\subsection{View Post Detail}

\textbf{a.~Class Diagram}

\kfig{0.92\textwidth}{images/sprint1/kvy/classDiagram/View Post Detail}{Class diagram of View Post Detail}{fig:cls-view-post-detail}

\textbf{b.~Sequence Diagram(s)}

\kfig{0.78\textwidth}{images/sprint1/kvy/sequences/PostDetail}{Sequence diagram of View Post Detail}{fig:seq-view-post-detail}

% ================= Create Post =================
\subsection{Create Post}

\textbf{a.~Class Diagram}

\kfig{0.92\textwidth}{images/sprint1/kvy/classDiagram/Create Post}{Class diagram of Create Post}{fig:cls-create-post}

\textbf{b.~Sequence Diagram(s)}

\kfig{0.78\textwidth}{images/sprint1/kvy/sequences/CreatePost}{Sequence diagram of Create Post}{fig:seq-create-post}

% ================= Update Post =================
\subsection{Update Post}

\textbf{a.~Class Diagram}

\kfig{0.92\textwidth}{images/sprint1/kvy/classDiagram/Update Post}{Class diagram of Update Post}{fig:cls-update-post}

\textbf{b.~Sequence Diagram(s)}

\kfig{0.78\textwidth}{images/sprint1/kvy/sequences/EditPost}{Sequence diagram of Update Post}{fig:seq-update-post}

% ================= Delete Post =================
\subsection{Delete Post}

\textbf{a.~Class Diagram}

\kfig{0.92\textwidth}{images/sprint1/kvy/classDiagram/Delete Post}{Class diagram of Delete Post}{fig:cls-delete-post}

\textbf{b.~Sequence Diagram(s)}

\kfig{0.78\textwidth}{images/sprint1/kvy/sequences/DeletePost}{Sequence diagram of Delete Post}{fig:seq-delete-post}

% ================= Like Post =================
\subsection{Like Post}

\textbf{a.~Class Diagram}

\kfig{0.92\textwidth}{images/sprint1/kvy/classDiagram/Like Post}{Class diagram of Like Post}{fig:cls-like-post}

\textbf{b.~Sequence Diagram(s)}

\kfig{0.78\textwidth}{images/sprint1/kvy/sequences/LikePost}{Sequence diagram of Like Post}{fig:seq-like-post}

% ================= Unlike Post =================
\subsection{Unlike Post}

\textbf{a.~Class Diagram}

\kfig{0.92\textwidth}{images/sprint1/kvy/classDiagram/Unlike Post}{Class diagram of Unlike Post}{fig:cls-unlike-post}

\textbf{b.~Sequence Diagram(s)}

\kfig{0.78\textwidth}{images/sprint1/kvy/sequences/UnlikePost}{Sequence diagram of Unlike Post}{fig:seq-unlike-post}

% ================= View Post Likes =================
\subsection{View Post Likes}

\textbf{a.~Class Diagram}

\kfig{0.92\textwidth}{images/sprint1/kvy/classDiagram/View Post Likes}{Class diagram of View Post Likes}{fig:cls-view-post-likes}

\textbf{b.~Sequence Diagram(s)}

\kfig{0.78\textwidth}{images/sprint1/kvy/sequences/ViewLikePostDetail}{Sequence diagram of View Post Likes}{fig:seq-view-post-likes}

% ================= View Post Comments =================
\subsection{View Post Comments}

\textbf{a.~Class Diagram}

\kfig{0.92\textwidth}{images/sprint1/kvy/classDiagram/View Post Comment}{Class diagram of View Post Comments}{fig:cls-view-post-comments}

\textbf{b.~Sequence Diagram(s)}

\kfig{0.78\textwidth}{images/sprint1/kvy/sequences/PostComments}{Sequence diagram of View Post Comments}{fig:seq-view-post-comments}

% ================= Create Comment =================
\subsection{Create Comment}

\textbf{a.~Class Diagram}

\kfig{0.92\textwidth}{images/sprint1/kvy/classDiagram/Create Comment}{Class diagram of Create Comment}{fig:cls-create-comment}

\textbf{b.~Sequence Diagram(s)}

\kfig{0.78\textwidth}{images/sprint1/kvy/sequences/AddComment}{Sequence diagram of Create Comment}{fig:seq-create-comment}

% ================= Update Comment =================
\subsection{Update Comment}

\textbf{a.~Class Diagram}

\kfig{0.92\textwidth}{images/sprint1/kvy/classDiagram/Update Comment}{Class diagram of Update Comment}{fig:cls-update-comment}

\textbf{b.~Sequence Diagram(s)}

\kfig{0.78\textwidth}{images/sprint1/kvy/sequences/EditComment}{Sequence diagram of Update Comment}{fig:seq-update-comment}

% ================= Delete Comment =================
\subsection{Delete Comment}

\textbf{a.~Class Diagram}

\kfig{0.92\textwidth}{images/sprint1/kvy/classDiagram/Delete Comment}{Class diagram of Delete Comment}{fig:cls-delete-comment}

\textbf{b.~Sequence Diagram(s)}

\kfig{0.78\textwidth}{images/sprint1/kvy/sequences/DeleteComment}{Sequence diagram of Delete Comment}{fig:seq-delete-comment}

% ================= Like Comment =================
\subsection{Like Comment}

\textbf{a.~Class Diagram}

\kfig{0.92\textwidth}{images/sprint1/kvy/classDiagram/Like Comment}{Class diagram of Like Comment}{fig:cls-like-comment}

\textbf{b.~Sequence Diagram(s)}

\kfig{0.78\textwidth}{images/sprint1/kvy/sequences/LikeComment}{Sequence diagram of Like Comment}{fig:seq-like-comment}

% ================= Unlike Comment =================
\subsection{Unlike Comment}

\textbf{a.~Class Diagram}

\kfig{0.92\textwidth}{images/sprint1/kvy/classDiagram/Unlike Comment}{Class diagram of Unlike Comment}{fig:cls-unlike-comment}

\textbf{b.~Sequence Diagram(s)}

\kfig{0.78\textwidth}{images/sprint1/kvy/sequences/UnlikeComment}{Sequence diagram of Unlike Comment}{fig:seq-unlike-comment}

% ================= Report Content =================
\subsection{Report Content}

\textbf{a.~Class Diagram}

\kfig{0.92\textwidth}{images/sprint1/kvy/classDiagram/Report Content}{Class diagram of Report Content}{fig:cls-report-content}

\textbf{b.~Sequence Diagram(s)}

\kfig{0.78\textwidth}{images/sprint1/kvy/sequences/ReportContent}{Sequence diagram of Report Content}{fig:seq-report-content}


% \input{sections/nhuy/sprint1}
% \input{sections/ntin/sprint1}
% \input{sections/ptai/sprint1}
% \input{sections/tthuan/sprint1}

% =================================================================
% FILE: sections/kvy/2/sprint2.tex -- Sprint 2 (Conversation)
% UC-27 Manage Conversation. Figures are the vector PDFs exported from
% class-diagrams/UC-27_Manage_Conversation_Class_Diagram.drawio and
% sequence-diagrams/UC-27.<nn>-*.drawio into images/sprint2/kvy/.
% =================================================================

% ================= Manage Conversation =================
\subsection{Manage Conversation}

Manage Conversation covers the UC-27 private messaging module: listing and
starting one-to-one conversations, opening a conversation, sending text
messages and files, marking messages as read, recalling a message, and
pinning or unpinning a conversation. Delivery is served both by the REST
endpoints and by the realtime gateway in the same namespace.

\subsubsection{Class Diagram}

\kfig{0.92\textwidth}{images/sprint2/kvy/classDiagram/ManageConversation}{Class diagram of Manage Conversation}{fig:kvy-cls-manage-conversation}

\subsubsection{UC-27.01 -- View Conversation List}
\kfig{0.78\textwidth}{images/sprint2/kvy/sequences/UC-27.01-view-conversation-list}{Sequence diagram of UC-27.01 View Conversation List}{fig:kvy-seq-uc-27-01}

\subsubsection{UC-27.02 -- Start Conversation}
\kfig{0.78\textwidth}{images/sprint2/kvy/sequences/UC-27.02-start-conversation}{Sequence diagram of UC-27.02 Start Conversation}{fig:kvy-seq-uc-27-02}

\subsubsection{UC-27.03 -- View Conversation Detail}
\kfig{0.78\textwidth}{images/sprint2/kvy/sequences/UC-27.03-view-conversation-detail}{Sequence diagram of UC-27.03 View Conversation Detail}{fig:kvy-seq-uc-27-03}

\subsubsection{UC-27.04 -- Send Message}
\kfig{0.78\textwidth}{images/sprint2/kvy/sequences/UC-27.04-send-message}{Sequence diagram of UC-27.04 Send Message}{fig:kvy-seq-uc-27-04}

\subsubsection{UC-27.05 -- Mark Message as Read}
\kfig{0.78\textwidth}{images/sprint2/kvy/sequences/UC-27.05-mark-message-as-read}{Sequence diagram of UC-27.05 Mark Message as Read}{fig:kvy-seq-uc-27-05}

\subsubsection{UC-27.06 -- Delete/Recall Message}
\kfig{0.78\textwidth}{images/sprint2/kvy/sequences/UC-27.06-delete-recall-message}{Sequence diagram of UC-27.06 Delete/Recall Message}{fig:kvy-seq-uc-27-06}

\subsubsection{UC-27.07 -- Pin/Unpin Conversation}
\kfig{0.78\textwidth}{images/sprint2/kvy/sequences/UC-27.07-pin-unpin-conversation}{Sequence diagram of UC-27.07 Pin/Unpin Conversation}{fig:kvy-seq-uc-27-07}

\subsubsection{UC-27.08 -- Send File}
\kfig{0.78\textwidth}{images/sprint2/kvy/sequences/UC-27.08-send-file}{Sequence diagram of UC-27.08 Send File}{fig:kvy-seq-uc-27-08}

% \input{sections/nhuy/sprint2}
% \input{sections/ntin/sprint2}
% \input{sections/ptai/sprint2}
% \input{sections/tthuan/sprint2}

% =================================================================
% FILE: sections/kvy/2/sprint3.tex -- Sprint 3 (Administration)
% UC-34 Manage Moderation, UC-35 Manage ID Verification and
% UC-36 Manage Forum Trash. Figures are the vector PDFs exported from
% class-diagrams/UC-3[456]_*.drawio and sequence-diagrams/UC-3[456].<nn>-*
% into images/sprint3/kvy/.
% =================================================================

% ================= Manage Moderation =================
\subsection{Manage Moderation}

Manage Moderation covers the UC-34 review workflow of the administration
console: listing the posts waiting for review, opening a flagged post and
approving or rejecting it. A review writes the moderation status, score and
reviewer back to the forum post.

\textbf{a.~Class Diagram}

\kfig{0.92\textwidth}{images/sprint3/kvy/classDiagram/ManageModeration}{Class diagram of Manage Moderation}{fig:kvy-cls-manage-moderation}

\textbf{b.~Sequence Diagram(s)}

\subsubsection{UC-34.01 -- View Moderation Queue}
\kfig{0.78\textwidth}{images/sprint3/kvy/sequences/UC-34.01-view-moderation-queue}{Sequence diagram of UC-34.01 View Moderation Queue}{fig:kvy-seq-uc-34-01}

\subsubsection{UC-34.02 -- View Moderation Item Detail}
\kfig{0.78\textwidth}{images/sprint3/kvy/sequences/UC-34.02-view-moderation-item-detail}{Sequence diagram of UC-34.02 View Moderation Item Detail}{fig:kvy-seq-uc-34-02}

\subsubsection{UC-34.03 -- Resolve Moderation Item}
\kfig{0.78\textwidth}{images/sprint3/kvy/sequences/UC-34.03-resolve-moderation-item}{Sequence diagram of UC-34.03 Resolve Moderation Item}{fig:kvy-seq-uc-34-03}

% ================= Manage ID Verification =================
\subsection{Manage ID Verification}

Manage ID Verification covers the UC-35 identity review workflow: listing the
submitted identity documents, opening a request with its uploaded file, and
approving or rejecting the request. A decision updates the verification status
and the reviewer of the identity document.

\textbf{a.~Class Diagram}

\kfig{0.92\textwidth}{images/sprint3/kvy/classDiagram/ManageIDVerification}{Class diagram of Manage ID Verification}{fig:kvy-cls-manage-id-verification}

\textbf{b.~Sequence Diagram(s)}

\subsubsection{UC-35.01 -- View ID Verification Request List}
\kfig{0.78\textwidth}{images/sprint3/kvy/sequences/UC-35.01-view-id-verification-request-list}{Sequence diagram of UC-35.01 View ID Verification Request List}{fig:kvy-seq-uc-35-01}

\subsubsection{UC-35.02 -- View ID Verification Request Detail}
\kfig{0.78\textwidth}{images/sprint3/kvy/sequences/UC-35.02-view-id-verification-request-detail}{Sequence diagram of UC-35.02 View ID Verification Request Detail}{fig:kvy-seq-uc-35-02}

\subsubsection{UC-35.03 -- Approve ID Verification}
\kfig{0.78\textwidth}{images/sprint3/kvy/sequences/UC-35.03-approve-id-verification}{Sequence diagram of UC-35.03 Approve ID Verification}{fig:kvy-seq-uc-35-03}

\subsubsection{UC-35.04 -- Reject ID Verification}
\kfig{0.78\textwidth}{images/sprint3/kvy/sequences/UC-35.04-reject-id-verification}{Sequence diagram of UC-35.04 Reject ID Verification}{fig:kvy-seq-uc-35-04}

% ================= Manage Forum Trash =================
\subsection{Manage Forum Trash}

Manage Forum Trash covers the UC-36 recovery workflow: listing the soft-deleted
posts and comments kept for review, opening one deleted item, and restoring it
so that it becomes visible in the forum again.

\textbf{a.~Class Diagram}

\kfig{0.92\textwidth}{images/sprint3/kvy/classDiagram/ManageForumTrash}{Class diagram of Manage Forum Trash}{fig:kvy-cls-manage-forum-trash}

\textbf{b.~Sequence Diagram(s)}

\subsubsection{UC-36.01 -- View Deleted Posts/Comments List}
\kfig{0.78\textwidth}{images/sprint3/kvy/sequences/UC-36.01-view-deleted-posts-comments-list}{Sequence diagram of UC-36.01 View Deleted Posts/Comments List}{fig:kvy-seq-uc-36-01}

\subsubsection{UC-36.02 -- View Deleted Posts/Comments Detail}
\kfig{0.78\textwidth}{images/sprint3/kvy/sequences/UC-36.02-view-deleted-posts-comments-detail}{Sequence diagram of UC-36.02 View Deleted Posts/Comments Detail}{fig:kvy-seq-uc-36-02}

\subsubsection{UC-36.03 -- Restore Deleted Posts/Comments}
\kfig{0.78\textwidth}{images/sprint3/kvy/sequences/UC-36.03-restore-deleted-posts-comments}{Sequence diagram of UC-36.03 Restore Deleted Posts/Comments}{fig:kvy-seq-uc-36-03}

% \input{sections/nhuy/sprint3}
% \input{sections/ntin/sprint3}
% \input{sections/ptai/sprint3}
% \input{sections/tthuan/sprint3}

% \input{sections/kvy/sprint4}
% \input{sections/nhuy/sprint4}
% \input{sections/ntin/sprint4}
% \input{sections/ptai/sprint4}
% \input{sections/tthuan/sprint4}

% \input{sections/nhuy/sprint5}
% \input{sections/ntin/sprint5}
% \input{sections/ptai/sprint5}
% \input{sections/tthuan/sprint5}

% ========================================================
% FILE: sections/4.Class_Specification.tex
% ========================================================
\section{Class Specification}

% ========================================================
% SPRINT 1
% ========================================================
% =================================================================
% FILE: sections/kvy/4/sprint1.tex -- Sprint 1 (Community Forum)
% Class specifications for the community-forum implementation.
% =================================================================

\subsection{ForumController Class}

\paragraph{ClassMethod}
\begin{ClassMethodTable}
01 & \code{ForumController(forumService: ForumService)} &
A method to initialize a new instance of the controller with the forum service.
\newline -- \code{forumService}: \code{ForumService} -- the service that owns category and post use cases.
\newline -- \textit{Returns:} \code{void} -- constructor, no return value. \\ \hline

02 & \code{getForumCategories()} &
A method to retrieve all active forum categories.
\newline -- \textit{Returns:} \code{Promise<ForumCategoryListResponseType>} -- the category list produced by the service.
\newline -- Route: \code{GET /forums/categories}; public endpoint. \\ \hline

03 & \code{getTopForumCategories(}\newline\code{query: { limit: number })} &
A method to retrieve the categories with the largest number of public posts.
\newline -- \code{query.limit}: \code{number} -- maximum number of categories, validated by \code{TopLimitSchema} with default 3 and range 1--10.
\newline -- \textit{Returns:} \code{Promise<ForumCategoryTopListResponseType>} -- ranked category list.
\newline -- Route: \code{GET /forums/categories/top}; public endpoint. \\ \hline

04 & \code{getTopActiveUsers(}\newline\code{query: { limit: number })} &
A method to retrieve the users with the most forum activity.
\newline -- \code{query.limit}: \code{number} -- maximum number of users, validated by \code{TopLimitSchema} with default 3 and range 1--10.
\newline -- \textit{Returns:} \code{Promise<ForumTopActiveUserListResponseType>} -- users with post and comment counts.
\newline -- Route: \code{GET /forums/users/top}; public endpoint. \\ \hline

05 & \code{getForumCategoryById(id: number)} &
A method to retrieve one forum category by its identifier.
\newline -- \code{id}: \code{number} -- category identifier parsed by \code{ParseIntPipe}.
\newline -- \textit{Returns:} \code{Promise<ForumCategoryDetailResponseType>} -- category detail with post count.
\newline -- Route: \code{GET /forums/categories/:id}; public endpoint. \\ \hline

06 & \code{getForumCategoryBySlug(slug: string)} &
A method to retrieve one forum category by its slug.
\newline -- \code{slug}: \code{string} -- unique category slug.
\newline -- \textit{Returns:} \code{Promise<ForumCategoryDetailResponseType>} -- category detail with post count.
\newline -- Route: \code{GET /forums/categories/by-slug/:slug}; public endpoint. \\ \hline

07 & \code{getForumPostLists(}\newline\code{filter: ForumPostFilterType)} &
A method to retrieve a paginated, filterable list of public posts.
\newline -- \code{filter}: \code{ForumPostFilterType} -- page, limit, optional categoryId, optional userId, and optional search text, validated by \code{ForumPostFilterDto}.
\newline -- \textit{Returns:} \code{Promise<ForumPostListResponseType>} -- posts plus pagination metadata.
\newline -- Route: \code{GET /forums/posts}; public endpoint. \\ \hline

08 & \code{getTopInteractedPosts(}\newline\code{query: { limit: number; categoryId?: number })} &
A method to retrieve the most interacted posts from the previous seven days.
\newline -- \code{query.limit}: \code{number} -- maximum number of posts, default 3.
\newline -- \code{query.categoryId}: \code{number | undefined} -- optional category filter.
\newline -- \textit{Returns:} \code{Promise<ForumTopPostListResponseType>} -- ranked posts with like and comment counts.
\newline -- Route: \code{GET /forums/posts/top}; public endpoint. \\ \hline

09 & \code{viewForumPostDetail(id: number)} &
A method to retrieve the public detail of one post.
\newline -- \code{id}: \code{number} -- post identifier.
\newline -- \textit{Returns:} \code{Promise<ViewForumPostDetailResponseType>} -- post, author summary, and category summary.
\newline -- Route: \code{GET /forums/posts/:id}; public endpoint. \\ \hline

10 & \code{getForumPostBySlug(slug: string)} &
A method to retrieve the public detail of one post by slug.
\newline -- \code{slug}: \code{string} -- unique post slug.
\newline -- \textit{Returns:} \code{Promise<ViewForumPostDetailResponseType>} -- post detail.
\newline -- Route: \code{GET /forums/posts/by-slug/:slug}; public endpoint. \\ \hline

11 & \code{createForumPost(}\newline\code{userId: number, body: CreateForumPostType)} &
A method to create a forum post for the authenticated user.
\newline -- \code{userId}: \code{number} -- author identifier injected by \code{@UserActive('userId')}.
\newline -- \code{body}: \code{CreateForumPostType} -- categoryId, title, content, and optional moderation result, validated by \code{CreateForumPostDto}.
\newline -- \textit{Returns:} \code{Promise<ForumPostType>} -- the created post.
\newline -- Route: \code{POST /forums/posts}; authenticated endpoint. \\ \hline

12 & \code{deleteForumPost(}\newline\code{userId: number, roleName: string,}\newline\code{id: number)} &
A method to soft-delete a post when the caller is its owner or an administrator.
\newline -- \code{userId}: \code{number} -- current user identifier.
\newline -- \code{roleName}: \code{string} -- current role name, used for the administrator override.
\newline -- \code{id}: \code{number} -- post identifier.
\newline -- \textit{Returns:} \code{Promise<ForumPostType>} -- the soft-deleted post.
\newline -- Route: \code{DELETE /forums/posts/:id}; authenticated endpoint. \\ \hline
\end{ClassMethodTable}

\paragraph{ClassAttribute}
\begin{ClassAttributeTable}
01 & \code{forumService} & \code{ForumService} -- private readonly service dependency. \\ \hline
\end{ClassAttributeTable}

\subsection{ForumService Class}

\paragraph{ClassMethod}
\begin{ClassMethodTable}
01 & \code{ForumService(forumRepository: ForumRepository)} &
A method to initialize the service with the forum repository.
\newline -- \code{forumRepository}: \code{ForumRepository} -- persistence dependency.
\newline -- \textit{Returns:} \code{void} -- constructor, no return value. \\ \hline

02 & \code{getForumCategories()} &
A method to load all active categories.
\newline -- \textit{Returns:} \code{Promise<ForumCategoryListResponseType>} -- category list with public post counts.
\newline -- Failure is mapped to \code{FailedToLoadForumCategoriesException}. \\ \hline

03 & \code{getTopForumCategories(limit: number = 3)} &
A method to load the categories with the most public posts.
\newline -- \code{limit}: \code{number} -- maximum number of categories; default 3.
\newline -- \textit{Returns:} \code{Promise<ForumCategoryTopListResponseType>} -- ranked categories.
\newline -- Failure is mapped to \code{FailedToLoadForumCategoriesException}. \\ \hline

04 & \code{getTopActiveUsers(limit: number = 5)} &
A method to load the most active users by post and comment counts.
\newline -- \code{limit}: \code{number} -- maximum number of users; default 5.
\newline -- \textit{Returns:} \code{Promise<ForumTopActiveUserListResponseType>} -- user activity summaries.
\newline -- The current implementation reuses \code{FailedToLoadForumCategoriesException} for this failure path. \\ \hline

05 & \code{getForumCategoryById(id: number)} &
A method to load one category by identifier.
\newline -- \code{id}: \code{number} -- category identifier.
\newline -- \textit{Returns:} \code{Promise<ForumCategoryDetailResponseType>} -- category detail.
\newline -- Prisma known errors are mapped to \code{FailedToLoadForumCategoriesException}. \\ \hline

06 & \code{getForumCategoryBySlug(slug: string)} &
A method to resolve a category slug and return its detail.
\newline -- \code{slug}: \code{string} -- category slug.
\newline -- \textit{Returns:} \code{Promise<ForumCategoryDetailResponseType>} -- category detail.
\newline -- Throws \code{ForumCategoryNotFoundException} when the slug does not exist. \\ \hline

07 & \code{getForumPostBySlug(slug: string)} &
A method to resolve a post slug and return its public detail.
\newline -- \code{slug}: \code{string} -- post slug.
\newline -- \textit{Returns:} \code{Promise<ViewForumPostDetailResponseType>} -- post detail.
\newline -- Throws \code{ForumPostNotFoundException} when the slug does not exist. \\ \hline

08 & \code{getForumPostLists(filter: ForumPostFilterType)} &
A method to build a paginated post list from the repository result.
\newline -- \code{filter}: \code{ForumPostFilterType} -- page, limit, category, author, and search filters.
\newline -- \textit{Returns:} \code{Promise<ForumPostListResponseType>} -- posts and pagination metadata, including \code{totalPages}.
\newline -- Prisma known errors are mapped to \code{FailedToLoadForumPostsException}. \\ \hline

09 & \code{createForumPost(}\newline\code{userId: number, body: CreateForumPostType)} &
A method to validate the optional category and create a post.
\newline -- \code{userId}: \code{number} -- author identifier.
\newline -- \code{body}: \code{CreateForumPostType} -- post data and optional AI moderation result.
\newline -- \textit{Returns:} \code{Promise<ForumPostType>} -- the created post.
\newline -- Throws \code{ForumCategoryNotFoundException} for an unknown category and \code{FailedToCreateForumPostException} for a persistence failure. \\ \hline

10 & \code{viewForumPostDetail(id: number)} &
A method to load the public detail of a post.
\newline -- \code{id}: \code{number} -- post identifier.
\newline -- \textit{Returns:} \code{Promise<ViewForumPostDetailResponseType>} -- post, author, and category summary.
\newline -- Prisma known errors are mapped to \code{FailedToViewForumPostException}. \\ \hline

11 & \code{updateForumPost(}\newline\code{id: number, userId: number,}\newline\code{body: UpdateForumPostType)} &
A method to update a post after checking existence and ownership.
\newline -- \code{id}: \code{number} -- post identifier.
\newline -- \code{userId}: \code{number} -- current user identifier.
\newline -- \code{body}: \code{UpdateForumPostType} -- updated category, title, and content.
\newline -- \textit{Returns:} \code{Promise<ForumPostType>} -- the updated post.
\newline -- Throws \code{ForumPostNotFoundException} or \code{ForumPostNotOwnedException}. The controller route is currently commented out. \\ \hline

12 & \code{deleteForumPost(}\newline\code{id: number, userId: number,}\newline\code{roleName: string)} &
A method to authorize and soft-delete a post.
\newline -- \code{id}: \code{number} -- post identifier.
\newline -- \code{userId}: \code{number} -- current user identifier.
\newline -- \code{roleName}: \code{string} -- current role name; \code{ADMIN} bypasses the owner check.
\newline -- \textit{Returns:} \code{Promise<ForumPostType>} -- the soft-deleted post.
\newline -- Throws \code{ForumPostNotFoundException}, \code{ForumPostNotOwnedException}, or \code{FailedToDeleteForumPostException}. \\ \hline

13 & \code{getTopInteractedPosts(}\newline\code{limit: number = 3,}\newline\code{categoryId?: number)} &
A method to return the most interacted posts from the last seven days.
\newline -- \code{limit}: \code{number} -- maximum number of posts; default 3.
\newline -- \code{categoryId}: \code{number | undefined} -- optional category restriction.
\newline -- \textit{Returns:} \code{Promise<ForumTopPostListResponseType>} -- ranked posts.
\newline -- Failure is mapped to \code{FailedToLoadForumPostsException}. \\ \hline
\end{ClassMethodTable}

\paragraph{ClassAttribute}
\begin{ClassAttributeTable}
01 & \code{forumRepository} & \code{ForumRepository} -- private readonly persistence dependency. \\ \hline
\end{ClassAttributeTable}

\subsection{ForumRepository Class}

\paragraph{ClassMethod}
\begin{ClassMethodTable}
01 & \code{ForumRepository(prisma: PrismaService)} &
A method to initialize the repository with the shared Prisma service.
\newline -- \code{prisma}: \code{PrismaService} -- pooled database client.
\newline -- \textit{Returns:} \code{void} -- constructor, no return value. \\ \hline

02 & \code{getForumCategories()} &
A method to query active categories ordered by name.
\newline -- \textit{Returns:} \code{Promise<Array<Pick<ForumCategoryType, 'id' | 'name' | 'slug' | 'description' | 'createdAt' | 'updatedAt' | 'postCount'>>>} -- category rows with public post counts.
\newline -- Only non-deleted posts with moderation status \code{APPROVED} are counted. \\ \hline

03 & \code{getTopForumCategories(limit: number = 3)} &
A method to query active categories ordered by public post count.
\newline -- \code{limit}: \code{number} -- maximum number of rows; default 3.
\newline -- \textit{Returns:} the same category projection as \code{getForumCategories()}, ranked by \code{posts._count} descending. \\ \hline

04 & \code{getTopActiveUsers(limit: number = 5)} &
A method to query non-deleted, non-banned users ordered by post count.
\newline -- \code{limit}: \code{number} -- maximum number of rows; default 5.
\newline -- \textit{Returns:} \code{Promise<ForumTopActiveUserType[]>} -- user id, display name, avatar, post count, and comment count.
\newline -- Comment counts include only non-deleted comments. \\ \hline

05 & \code{findForumCategoryById(id: number)} &
A method to check whether an active category exists.
\newline -- \code{id}: \code{number} -- category identifier.
\newline -- \textit{Returns:} \code{Promise<Pick<ForumCategoryType, 'id'> | null>} -- the id when found, otherwise \code{null}. \\ \hline

06 & \code{getForumCategoryById(id: number)} &
A method to load one active category with its post count.
\newline -- \code{id}: \code{number} -- category identifier.
\newline -- \textit{Returns:} category projection with \code{postCount}.
\newline -- Throws \code{ForumCategoryNotFoundException} when no active row exists. \\ \hline

07 & \code{findForumCategoryBySlug(slug: string)} &
A method to resolve an active category by slug.
\newline -- \code{slug}: \code{string} -- category slug.
\newline -- \textit{Returns:} \code{Promise<Pick<ForumCategoryType, 'id'> | null>} -- the id when found, otherwise \code{null}. \\ \hline

08 & \code{getForumPostLists(filter: ForumPostFilterType)} &
A method to query a page of public posts and the matching total in one transaction.
\newline -- \code{filter}: \code{ForumPostFilterType} -- pagination, category, author, and case-insensitive title search.
\newline -- \textit{Returns:} \code{Promise<{ posts: ForumPostWithUserType[]; total: number }>} -- mapped post rows and total count.
\newline -- Visibility requires \code{deletedAt = null} and moderation status \code{APPROVED}. \\ \hline

09 & \code{createForumPost(}\newline\code{categoryId: number | null, userId: number,}\newline\code{title: string, content: string,}\newline\code{moderation?: ForumModerationType)} &
A method to insert a post and generate its slug.
\newline -- \code{categoryId}: \code{number | null} -- optional category.
\newline -- \code{userId}: \code{number} -- author identifier.
\newline -- \code{title}, \code{content}: \code{string} -- post text.
\newline -- \code{moderation}: \code{ForumModerationType | undefined} -- optional AI result; status defaults to \code{PENDING}.
\newline -- \textit{Returns:} \code{Promise<ForumPostType>} -- the inserted post. \\ \hline

10 & \code{viewForumPostDetail(id: number)} &
A method to query one public post with author and category summaries.
\newline -- \code{id}: \code{number} -- post identifier.
\newline -- \textit{Returns:} \code{Promise<ViewForumPostDetailResponseType>} -- post detail.
\newline -- Throws \code{ForumPostNotFoundException} when the post is missing, deleted, or not approved. \\ \hline

11 & \code{findForumPostById(id: number)} &
A method to load a non-deleted post for ownership checks.
\newline -- \code{id}: \code{number} -- post identifier.
\newline -- \textit{Returns:} \code{Promise<ForumPostType | null>} -- the post or \code{null}. \\ \hline

12 & \code{findForumPostBySlug(slug: string)} &
A method to resolve a public post slug.
\newline -- \code{slug}: \code{string} -- post slug.
\newline -- \textit{Returns:} \code{Promise<Pick<ForumPostType, 'id'> | null>} -- the id when found, otherwise \code{null}. \\ \hline

13 & \code{updateForumPost(}\newline\code{id: number, data: UpdateForumPostType)} &
A method to update a non-deleted post and regenerate its slug from the title.
\newline -- \code{id}: \code{number} -- post identifier.
\newline -- \code{data}: \code{UpdateForumPostType} -- category, title, and content.
\newline -- \textit{Returns:} \code{Promise<ForumPostType>} -- the updated post. \\ \hline

14 & \code{softDeleteForumPost(id: number)} &
A method to set \code{deletedAt} on a non-deleted post.
\newline -- \code{id}: \code{number} -- post identifier.
\newline -- \textit{Returns:} \code{Promise<ForumPostType>} -- the soft-deleted post. \\ \hline

15 & \code{getTopInteractedPosts(}\newline\code{limit: number = 3,}\newline\code{categoryId?: number)} &
A method to query posts created during the previous seven days and rank them by interactions.
\newline -- \code{limit}: \code{number} -- maximum number of rows; default 3.
\newline -- \code{categoryId}: \code{number | undefined} -- optional category restriction.
\newline -- \textit{Returns:} \code{Promise<ForumTopPostType[]>} -- posts with like count, comment count, and \code{interactionScore}.
\newline -- Ordering is by likes, then comments, then creation time, all descending. \\ \hline
\end{ClassMethodTable}

\paragraph{ClassAttribute}
\begin{ClassAttributeTable}
01 & \code{prisma} & \code{PrismaService} -- private readonly database client. \\ \hline
\end{ClassAttributeTable}

\subsubsection*{Supporting DTO classes for category and post management}

\begin{SupportingTable}{4.0cm}{7.2cm}{4.6cm}
\code{ForumCategoryListResponseDto} & Array of \code{ForumCategoryType}: \code{id}, \code{name}, \code{slug}, \code{description}, \code{createdAt}, \code{updatedAt}, \code{postCount}. & Response of \code{GET /forums/categories}. \\ \hline
\code{ForumCategoryDetailResponseDto} & \code{id: number}, \code{name: string}, \code{slug: string}, \code{description: string | null}, \code{createdAt: Date}, \code{updatedAt: Date}, \code{postCount: number}. & Detail response for category id and slug routes. \\ \hline
\code{ForumCategoryTopListResponseDto} & Array of \code{ForumCategoryType}. & Response of the top-category route. \\ \hline
\code{ForumTopActiveUserListResponseDto} & Array of \code{ForumTopActiveUserType}: \code{id}, \code{displayName}, \code{avatarUrl}, \code{postCount}, \code{commentCount}. & Response of the top-active-user route. \\ \hline
\code{ForumPostListResponseDto} & \code{posts: ForumPostWithUserType[]}, \code{pagination: PaginationType}. & Paginated post-list response. \\ \hline
\code{ForumPostFilterDto} & \code{page: number}, \code{limit: number}, \code{categoryId?: number}, \code{userId?: number}, \code{search?: string}. & Query validation for the post list. \\ \hline
\code{CreateForumPostDto} & \code{categoryId?: number}, \code{title: string}, \code{content: string}, \code{moderation?: ForumModerationType}. & Request body for post creation. \\ \hline
\code{CreateForumPostResponseDto} & \code{ForumPostType}: \code{id}, \code{categoryId}, \code{userId}, \code{title}, \code{slug}, \code{content}, \code{moderationStatus}, \code{moderationScore}, \code{moderationCategories}, \code{createdAt}, \code{updatedAt}. & Response of post creation. \\ \hline
\code{ViewForumPostDetailResponseDto} & \code{ForumPostType} fields plus \code{category: { id, name, slug } | null} and \code{user: { id, profile }}. & Public post-detail response. \\ \hline
\code{UpdateForumPostDto} & \code{categoryId?: number}, \code{title: string}, \code{content: string}. & Request body for the implemented update-post method. \\ \hline
\code{UpdateForumPostResponseDto} & \code{ForumPostType}. & Response of the implemented update-post method. \\ \hline
\code{ForumTopPostListResponseDto} & Array of \code{ForumTopPostType}: post fields plus \code{likeCount}, \code{commentCount}, \code{interactionScore}, \code{user}, and \code{category}. & Response of the top-post route. \\ \hline
\end{SupportingTable}

\subsection{ForumLikeController Class}

\paragraph{ClassMethod}
\begin{ClassMethodTable}
01 & \code{ForumLikeController(forumLikeService: ForumLikeService)} &
A method to initialize the controller with the post-like service.
\newline -- \code{forumLikeService}: \code{ForumLikeService} -- service dependency.
\newline -- \textit{Returns:} \code{void} -- constructor, no return value. \\ \hline

02 & \code{toggleLikePost(userId: number, id: number)} &
A method to toggle the current user's like on a post.
\newline -- \code{userId}: \code{number} -- current user identifier.
\newline -- \code{id}: \code{number} -- post identifier.
\newline -- \textit{Returns:} \code{Promise<ToggleLikeResponseType>} -- \code{liked: true} after a like, \code{liked: false} after an unlike.
\newline -- Route: \code{POST /forums/posts/:id/like}; authenticated endpoint. \\ \hline

03 & \code{viewForumLikeDetail(postId: number)} &
A method to list the users who liked a post.
\newline -- \code{postId}: \code{number} -- post identifier.
\newline -- \textit{Returns:} \code{Promise<ForumLikeDetailResponseType>} -- like rows with user summaries.
\newline -- Route: \code{GET /forums/posts/:id/like}; public endpoint. \\ \hline
\end{ClassMethodTable}

\paragraph{ClassAttribute}
\begin{ClassAttributeTable}
01 & \code{forumLikeService} & \code{ForumLikeService} -- private readonly service dependency. \\ \hline
\end{ClassAttributeTable}

\subsection{ForumLikeService Class}

\paragraph{ClassMethod}
\begin{ClassMethodTable}
01 & \code{ForumLikeService(forumLikeRepository: ForumLikeRepository)} &
A method to initialize the service with the post-like repository.
\newline -- \code{forumLikeRepository}: \code{ForumLikeRepository} -- persistence dependency.
\newline -- \textit{Returns:} \code{void} -- constructor, no return value. \\ \hline

02 & \code{toggleLikePost(userId: number, postId: number)} &
A method to decide whether a like must be created or removed.
\newline -- \code{userId}: \code{number} -- current user identifier.
\newline -- \code{postId}: \code{number} -- post identifier.
\newline -- \textit{Returns:} \code{Promise<ToggleLikeResponseType>} -- the new like state.
\newline -- Throws \code{ForumPostNotFoundException} when the post is missing; other failures become \code{FailedToLikeForumPostException}. \\ \hline

03 & \code{viewForumLikeDetail(postId: number)} &
A method to load the like rows of a post.
\newline -- \code{postId}: \code{number} -- post identifier.
\newline -- \textit{Returns:} \code{Promise<ForumLikeDetailResponseType>} -- like rows with user summaries.
\newline -- Failure is mapped to \code{FailedToViewForumLikeDetailException}. \\ \hline
\end{ClassMethodTable}

\paragraph{ClassAttribute}
\begin{ClassAttributeTable}
01 & \code{forumLikeRepository} & \code{ForumLikeRepository} -- private readonly persistence dependency. \\ \hline
\end{ClassAttributeTable}

\subsection{ForumLikeRepository Class}

\paragraph{ClassMethod}
\begin{ClassMethodTable}
01 & \code{ForumLikeRepository(prisma: PrismaService)} &
A method to initialize the repository with the shared Prisma service.
\newline -- \code{prisma}: \code{PrismaService} -- database client.
\newline -- \textit{Returns:} \code{void} -- constructor, no return value. \\ \hline

02 & \code{findPostById(postId: number)} &
A method to check that a non-deleted post exists.
\newline -- \code{postId}: \code{number} -- post identifier.
\newline -- \textit{Returns:} \code{Promise<{ id: number } | null>} -- the id when found, otherwise \code{null}. \\ \hline

03 & \code{findLikeByUserAndPost(}\newline\code{userId: number, postId: number)} &
A method to find a user's reaction row for a post.
\newline -- \code{userId}: \code{number} -- user identifier.
\newline -- \code{postId}: \code{number} -- post identifier.
\newline -- \textit{Returns:} \code{Promise<{ id: number; commentId: number | null } | null>} -- reaction id and optional comment link. \\ \hline

04 & \code{findCommentOnlyLike(userId: number)} &
A method to find a reaction row that currently links only to a comment.
\newline -- \code{userId}: \code{number} -- user identifier.
\newline -- \textit{Returns:} \code{Promise<{ id: number } | null>} -- the reusable reaction row, otherwise \code{null}. \\ \hline

05 & \code{createPostLike(userId: number, postId: number)} &
A method to create a post reaction, reusing a comment-only row when possible.
\newline -- \code{userId}: \code{number} -- user identifier.
\newline -- \code{postId}: \code{number} -- post identifier.
\newline -- \textit{Returns:} the created or updated \code{ForumLike} row. \\ \hline

06 & \code{deletePostLike(userId: number, postId: number)} &
A method to remove a post reaction while preserving a comment reaction on the same row.
\newline -- \code{userId}: \code{number} -- user identifier.
\newline -- \code{postId}: \code{number} -- post identifier.
\newline -- \textit{Returns:} the updated or deleted \code{ForumLike} row, or \code{undefined} when no row exists. \\ \hline

07 & \code{viewForumLikeDetail(postId: number)} &
A method to query all reaction rows for a post.
\newline -- \code{postId}: \code{number} -- post identifier.
\newline -- \textit{Returns:} \code{Promise<ForumLikeDetailResponseType>} -- reaction rows with user id, profile display name, and avatar. \\ \hline
\end{ClassMethodTable}

\paragraph{ClassAttribute}
\begin{ClassAttributeTable}
01 & \code{prisma} & \code{PrismaService} -- private readonly database client. \\ \hline
\end{ClassAttributeTable}

\subsubsection*{Supporting DTO classes for post reactions}

\begin{SupportingTable}{4.4cm}{7.4cm}{4.0cm}
\code{ToggleLikeResponseDto} & \code{liked: boolean}. & Response of the toggle-like route. \\ \hline
\code{ViewForumLikeDetailResponseDto} & Array of \code{ForumLikeType} with \code{user}: \code{id}, \code{userId}, \code{postId}, \code{commentId}, \code{reactionType}, \code{createdAt}, and user profile summary. & Response of the post-like detail route. \\ \hline
\end{SupportingTable}

\subsection{ForumCommentController Class}

\paragraph{ClassMethod}
\begin{ClassMethodTable}
01 & \code{ForumCommentController(}\newline\code{forumCommentService: ForumCommentService)} &
A method to initialize the controller with the comment service.
\newline -- \code{forumCommentService}: \code{ForumCommentService} -- service dependency.
\newline -- \textit{Returns:} \code{void} -- constructor, no return value. \\ \hline

02 & \code{getForumCommentLists(}\newline\code{postId: number,}\newline\code{filter: ForumCommentFilterType,}\newline\code{userId: number)} &
A method to list the comments of a post with pagination.
\newline -- \code{postId}: \code{number} -- parent post identifier.
\newline -- \code{filter}: \code{ForumCommentFilterType} -- page and limit.
\newline -- \code{userId}: \code{number} -- current user, used to compute \code{likedByMe}.
\newline -- \textit{Returns:} \code{Promise<ForumCommentListResponseType>} -- comments and pagination.
\newline -- Route: \code{GET /forums/posts/:postId/comments}; authenticated endpoint. \\ \hline

03 & \code{createForumComment(}\newline\code{userId: number, postId: number,}\newline\code{body: CreateForumCommentType)} &
A method to create a comment on a post.
\newline -- \code{userId}: \code{number} -- author identifier.
\newline -- \code{postId}: \code{number} -- parent post identifier.
\newline -- \code{body}: \code{CreateForumCommentType} -- comment content.
\newline -- \textit{Returns:} \code{Promise<ForumCommentType>} -- the created comment.
\newline -- Route: \code{POST /forums/posts/:postId/comments}; authenticated endpoint. \\ \hline

04 & \code{updateForumComment(}\newline\code{userId: number, postId: number,}\newline\code{id: number, body: EditForumCommentType)} &
A method to edit a comment owned by the current user.
\newline -- \code{userId}: \code{number} -- current user identifier.
\newline -- \code{postId}: \code{number} -- parent post identifier.
\newline -- \code{id}: \code{number} -- comment identifier.
\newline -- \code{body}: \code{EditForumCommentType} -- new content.
\newline -- \textit{Returns:} \code{Promise<ForumCommentType>} -- the updated comment.
\newline -- Route: \code{PATCH /forums/posts/:postId/comments/:id}; authenticated endpoint. \\ \hline

05 & \code{deleteForumComment(}\newline\code{userId: number, roleName: string,}\newline\code{postId: number, id: number)} &
A method to soft-delete a comment owned by the user or by an administrator.
\newline -- \code{userId}: \code{number} -- current user identifier.
\newline -- \code{roleName}: \code{string} -- current role name.
\newline -- \code{postId}: \code{number} -- parent post identifier.
\newline -- \code{id}: \code{number} -- comment identifier.
\newline -- \textit{Returns:} \code{Promise<ForumCommentType>} -- the soft-deleted comment.
\newline -- Route: \code{DELETE /forums/posts/:postId/comments/:id}; authenticated endpoint. \\ \hline

06 & \code{toggleLikeComment(}\newline\code{userId: number, postId: number,}\newline\code{commentId: number)} &
A method to toggle the current user's like on a comment.
\newline -- \code{userId}: \code{number} -- current user identifier.
\newline -- \code{postId}: \code{number} -- parent post identifier.
\newline -- \code{commentId}: \code{number} -- comment identifier.
\newline -- \textit{Returns:} \code{Promise<ToggleLikeCommentType>} -- the new like state.
\newline -- Route: \code{POST /forums/posts/:postId/comments/:commentId/like}; authenticated endpoint. \\ \hline
\end{ClassMethodTable}

\paragraph{ClassAttribute}
\begin{ClassAttributeTable}
01 & \code{forumCommentService} & \code{ForumCommentService} -- private readonly service dependency. \\ \hline
\end{ClassAttributeTable}

\subsection{ForumCommentService Class}

\paragraph{ClassMethod}
\begin{ClassMethodTable}
01 & \code{ForumCommentService(}\newline\code{forumCommentRepository: ForumCommentRepository)} &
A method to initialize the service with the comment repository.
\newline -- \code{forumCommentRepository}: \code{ForumCommentRepository} -- persistence dependency.
\newline -- \textit{Returns:} \code{void} -- constructor, no return value. \\ \hline

02 & \code{createForumComment(}\newline\code{postId: number, userId: number,}\newline\code{body: CreateForumCommentType)} &
A method to verify the parent post and create a comment.
\newline -- \code{postId}: \code{number} -- parent post identifier.
\newline -- \code{userId}: \code{number} -- author identifier.
\newline -- \code{body}: \code{CreateForumCommentType} -- comment content.
\newline -- \textit{Returns:} \code{Promise<ForumCommentType>} -- the created comment.
\newline -- Throws \code{ForumPostNotFoundException} or \code{FailedToCreateForumCommentException}. \\ \hline

03 & \code{updateForumComment(}\newline\code{postId: number, id: number,}\newline\code{userId: number, body: EditForumCommentType)} &
A method to verify comment existence, parent post, and ownership before editing.
\newline -- \code{postId}: \code{number} -- parent post identifier.
\newline -- \code{id}: \code{number} -- comment identifier.
\newline -- \code{userId}: \code{number} -- current user identifier.
\newline -- \code{body}: \code{EditForumCommentType} -- new content.
\newline -- \textit{Returns:} \code{Promise<ForumCommentType>} -- the updated comment.
\newline -- Throws \code{ForumCommentNotFoundException} or \code{ForumCommentNotOwnedException}. \\ \hline

04 & \code{deleteForumComment(}\newline\code{postId: number, id: number,}\newline\code{userId: number, roleName: string)} &
A method to authorize and soft-delete a comment.
\newline -- \code{postId}: \code{number} -- parent post identifier.
\newline -- \code{id}: \code{number} -- comment identifier.
\newline -- \code{userId}: \code{number} -- current user identifier.
\newline -- \code{roleName}: \code{string} -- \code{ADMIN} bypasses the owner check.
\newline -- \textit{Returns:} \code{Promise<ForumCommentType>} -- the soft-deleted comment.
\newline -- Throws \code{ForumCommentNotFoundException}, \code{ForumCommentNotOwnedException}, or \code{FailedToDeleteForumCommentException}. \\ \hline

05 & \code{toggleLikeComment(}\newline\code{userId: number, postId: number,}\newline\code{commentId: number)} &
A method to create or remove a comment reaction.
\newline -- \code{userId}: \code{number} -- current user identifier.
\newline -- \code{postId}: \code{number} -- parent post identifier.
\newline -- \code{commentId}: \code{number} -- comment identifier.
\newline -- \textit{Returns:} \code{Promise<ToggleLikeCommentType>} -- the new like state.
\newline -- Throws \code{ForumCommentNotFoundException} or \code{FailedToLikeForumCommentException}. \\ \hline

06 & \code{getForumCommentLists(filter: ForumCommentFilterType)} &
A method to build a paginated comment list.
\newline -- \code{filter}: \code{ForumCommentFilterType} -- page, limit, postId, and optional userId.
\newline -- \textit{Returns:} \code{Promise<ForumCommentListResponseType>} -- comments and pagination metadata.
\newline -- Prisma known errors are mapped to \code{FailedToLoadForumCommentsException}. \\ \hline
\end{ClassMethodTable}

\paragraph{ClassAttribute}
\begin{ClassAttributeTable}
01 & \code{forumCommentRepository} & \code{ForumCommentRepository} -- private readonly persistence dependency. \\ \hline
\end{ClassAttributeTable}

\subsection{ForumCommentRepository Class}

\paragraph{ClassMethod}
\begin{ClassMethodTable}
01 & \code{ForumCommentRepository(prisma: PrismaService)} &
A method to initialize the repository with the shared Prisma service.
\newline -- \code{prisma}: \code{PrismaService} -- database client.
\newline -- \textit{Returns:} \code{void} -- constructor, no return value. \\ \hline

02 & \code{findPostById(postId: number)} &
A method to check that a non-deleted post exists.
\newline -- \code{postId}: \code{number} -- post identifier.
\newline -- \textit{Returns:} \code{Promise<{ id: number } | null>} -- the id when found, otherwise \code{null}. \\ \hline

03 & \code{findCommentById(id: number)} &
A method to load a non-deleted comment with its parent post and author summary.
\newline -- \code{id}: \code{number} -- comment identifier.
\newline -- \textit{Returns:} the comment projection or \code{null}. \\ \hline

04 & \code{createForumComment(}\newline\code{postId: number, userId: number,}\newline\code{content: string)} &
A method to insert a comment and return the response shape.
\newline -- \code{postId}: \code{number} -- parent post identifier.
\newline -- \code{userId}: \code{number} -- author identifier.
\newline -- \code{content}: \code{string} -- comment text.
\newline -- \textit{Returns:} \code{Promise<ForumCommentType>} -- the comment with \code{likeCount = 0} and \code{likedByMe = false}. \\ \hline

05 & \code{updateForumComment(}\newline\code{id: number, data: EditForumCommentType)} &
A method to update the content of a non-deleted comment.
\newline -- \code{id}: \code{number} -- comment identifier.
\newline -- \code{data}: \code{EditForumCommentType} -- new content.
\newline -- \textit{Returns:} \code{Promise<ForumCommentType>} -- the updated comment. \\ \hline

06 & \code{softDeleteForumComment(id: number)} &
A method to set \code{deletedAt} on a non-deleted comment.
\newline -- \code{id}: \code{number} -- comment identifier.
\newline -- \textit{Returns:} \code{Promise<ForumCommentType>} -- the soft-deleted comment. \\ \hline

07 & \code{getForumCommentLists(}\newline\code{filter: ForumCommentFilterType)} &
A method to query a page of comments and the matching total in one transaction.
\newline -- \code{filter}: \code{ForumCommentFilterType} -- page, limit, postId, and optional current user.
\newline -- \textit{Returns:} \code{Promise<{ comments: ForumCommentType[]; total: number }>} -- mapped comments and total.
\newline -- \code{likedByMe} is computed from the optional user reaction. \\ \hline

08 & \code{findLikeByUserAndComment(}\newline\code{userId: number, commentId: number)} &
A method to find a user's comment reaction.
\newline -- \code{userId}: \code{number} -- user identifier.
\newline -- \code{commentId}: \code{number} -- comment identifier.
\newline -- \textit{Returns:} \code{Promise<{ id: number; postId: number | null } | null>} -- reaction id and optional post link. \\ \hline

09 & \code{findLikeByUserAndPost(}\newline\code{userId: number, postId: number)} &
A method to find a user's post reaction for reuse by the comment-like flow.
\newline -- \code{userId}: \code{number} -- user identifier.
\newline -- \code{postId}: \code{number} -- post identifier.
\newline -- \textit{Returns:} \code{Promise<{ id: number; commentId: number | null } | null>} -- reaction id and optional comment link. \\ \hline

10 & \code{createCommentLike(}\newline\code{userId: number, postId: number,}\newline\code{commentId: number)} &
A method to create a comment reaction, reusing a post-only row when possible.
\newline -- \code{userId}: \code{number} -- user identifier.
\newline -- \code{postId}: \code{number} -- parent post identifier.
\newline -- \code{commentId}: \code{number} -- comment identifier.
\newline -- \textit{Returns:} the created or updated \code{ForumLike} row. \\ \hline

11 & \code{deleteCommentLike(}\newline\code{userId: number, commentId: number)} &
A method to remove a comment reaction while preserving a post reaction on the same row.
\newline -- \code{userId}: \code{number} -- user identifier.
\newline -- \code{commentId}: \code{number} -- comment identifier.
\newline -- \textit{Returns:} the updated or deleted \code{ForumLike} row, or \code{undefined} when no row exists. \\ \hline
\end{ClassMethodTable}

\paragraph{ClassAttribute}
\begin{ClassAttributeTable}
01 & \code{prisma} & \code{PrismaService} -- private readonly database client. \\ \hline
\end{ClassAttributeTable}

\subsubsection*{Supporting DTO classes for comments}

\begin{SupportingTable}{4.2cm}{7.4cm}{4.2cm}
\code{ForumCommentListResponseDto} & \code{comments: ForumCommentType[]}, \code{pagination: PaginationType}. & Paginated comment-list response. \\ \hline
\code{ForumCommentFilterDto} & \code{page: number}, \code{limit: number}. & Query validation for comment listing. \\ \hline
\code{CreateForumCommentDto} & \code{content: string}. & Comment-creation body. \\ \hline
\code{CreateForumCommentResponseDto} & \code{ForumCommentType}: \code{id}, \code{postId}, \code{user}, \code{content}, \code{createdAt}, \code{updatedAt}, \code{likeCount}, \code{likedByMe}. & Comment-creation response. \\ \hline
\code{EditForumCommentDto} & \code{content: string}. & Comment-update body. \\ \hline
\code{EditForumCommentResponseDto} & \code{ForumCommentType}. & Comment-update response. \\ \hline
\code{DeleteForumCommentResponseDto} & \code{ForumCommentType}. & Soft-delete response. \\ \hline
\code{ToggleLikeCommentResponseDto} & \code{liked: boolean}. & Comment-like toggle response. \\ \hline
\end{SupportingTable}

\subsection{ForumReportController Class}

\paragraph{ClassMethod}
\begin{ClassMethodTable}
01 & \code{ForumReportController(}\newline\code{forumReportService: ForumReportService)} &
A method to initialize the controller with the report service.
\newline -- \code{forumReportService}: \code{ForumReportService} -- service dependency.
\newline -- \textit{Returns:} \code{void} -- constructor, no return value. \\ \hline

02 & \code{createPostReport(}\newline\code{userId: number, postId: number,}\newline\code{body: CreateForumReportType)} &
A method to report a post.
\newline -- \code{userId}: \code{number} -- reporter identifier.
\newline -- \code{postId}: \code{number} -- reported post identifier.
\newline -- \code{body}: \code{CreateForumReportType} -- report reason.
\newline -- \textit{Returns:} \code{Promise<CreateForumReportResponseType>} -- the created report.
\newline -- Route: \code{POST /forums/posts/:postId/reports}; authenticated endpoint. \\ \hline

03 & \code{createCommentReport(}\newline\code{userId: number, postId: number,}\newline\code{commentId: number,}\newline\code{body: CreateForumReportType)} &
A method to report a comment.
\newline -- \code{userId}: \code{number} -- reporter identifier.
\newline -- \code{postId}: \code{number} -- parent post identifier.
\newline -- \code{commentId}: \code{number} -- reported comment identifier.
\newline -- \code{body}: \code{CreateForumReportType} -- report reason.
\newline -- \textit{Returns:} \code{Promise<CreateForumReportResponseType>} -- the created report.
\newline -- Route: \code{POST /forums/posts/:postId/comments/:commentId/reports}; authenticated endpoint. \\ \hline

04 & \code{checkPostReported(}\newline\code{userId: number, postId: number)} &
A method to check whether the current user has reported a post.
\newline -- \code{userId}: \code{number} -- current user identifier.
\newline -- \code{postId}: \code{number} -- post identifier.
\newline -- \textit{Returns:} \code{Promise<{ reported: boolean }>} -- true when a matching report exists.
\newline -- Route: \code{GET /forums/posts/:postId/is-reported}; authenticated endpoint. \\ \hline

05 & \code{checkCommentReported(}\newline\code{userId: number, postId: number,}\newline\code{commentId: number)} &
A method to check whether the current user has reported a comment.
\newline -- \code{userId}: \code{number} -- current user identifier.
\newline -- \code{postId}: \code{number} -- parent post identifier.
\newline -- \code{commentId}: \code{number} -- comment identifier.
\newline -- \textit{Returns:} \code{Promise<{ reported: boolean }>} -- true when a matching report exists.
\newline -- Route: \code{GET /forums/posts/:postId/comments/:commentId/is-reported}; authenticated endpoint. \\ \hline

06 & \code{getForumReportLists(}\newline\code{userId: number, roleName: string,}\newline\code{page: number, limit: number,}\newline\code{status: string, search: string)} &
A method to list the report queue for administrators.
\newline -- \code{userId}: \code{number} -- current user identifier.
\newline -- \code{roleName}: \code{string} -- role checked by the service.
\newline -- \code{page}, \code{limit}: \code{number} -- pagination, defaulted to 1 and 10.
\newline -- \code{status}, \code{search}: \code{string} -- optional status and reason filters.
\newline -- \textit{Returns:} \code{Promise<ForumReportListResponseType>} -- reports and pagination.
\newline -- Route: \code{GET /forums/reports}; administrator endpoint. \\ \hline

07 & \code{updateReportStatus(}\newline\code{userId: number, roleName: string,}\newline\code{reportId: number,}\newline\code{body: UpdateReportStatusType)} &
A method to change the status of a report.
\newline -- \code{userId}: \code{number} -- administrator identifier recorded on the report.
\newline -- \code{roleName}: \code{string} -- role checked by the service.
\newline -- \code{reportId}: \code{number} -- report identifier.
\newline -- \code{body}: \code{UpdateReportStatusType} -- target status.
\newline -- \textit{Returns:} \code{Promise<UpdateReportStatusResponseType>} -- the updated report.
\newline -- Route: \code{PATCH /forums/reports/:reportId/status}; administrator endpoint. \\ \hline
\end{ClassMethodTable}

\paragraph{ClassAttribute}
\begin{ClassAttributeTable}
01 & \code{forumReportService} & \code{ForumReportService} -- private readonly service dependency. \\ \hline
\end{ClassAttributeTable}

\subsection{ForumReportService Class}

\paragraph{ClassMethod}
\begin{ClassMethodTable}
01 & \code{ForumReportService(}\newline\code{forumReportRepository: ForumReportRepository)} &
A method to initialize the service with the report repository.
\newline -- \code{forumReportRepository}: \code{ForumReportRepository} -- persistence dependency.
\newline -- \textit{Returns:} \code{void} -- constructor, no return value. \\ \hline

02 & \code{createPostReport(}\newline\code{reporterId: number, postId: number,}\newline\code{reason: string)} &
A method to validate a post and create its first report by the user.
\newline -- \code{reporterId}: \code{number} -- reporter identifier.
\newline -- \code{postId}: \code{number} -- reported post identifier.
\newline -- \code{reason}: \code{string} -- report reason.
\newline -- \textit{Returns:} \code{Promise<CreateForumReportResponseType>} -- the created report.
\newline -- Throws \code{ForumPostNotFoundException}, \code{ForumReportAlreadyExistsException}, or \code{FailedToCreateForumReportException}. \\ \hline

03 & \code{createCommentReport(}\newline\code{reporterId: number, postId: number,}\newline\code{commentId: number, reason: string)} &
A method to validate a comment and create its first report by the user.
\newline -- \code{reporterId}: \code{number} -- reporter identifier.
\newline -- \code{postId}: \code{number} -- parent post identifier.
\newline -- \code{commentId}: \code{number} -- reported comment identifier.
\newline -- \code{reason}: \code{string} -- report reason.
\newline -- \textit{Returns:} \code{Promise<CreateForumReportResponseType>} -- the created report.
\newline -- Throws \code{ForumCommentNotFoundException}, \code{ForumReportAlreadyExistsException}, or \code{FailedToCreateForumReportException}. \\ \hline

04 & \code{checkReportedByUser(}\newline\code{reporterId: number, postId: number | null,}\newline\code{commentId: number | null)} &
A method to find an existing report for the current user and target.
\newline -- \code{reporterId}: \code{number} -- reporter identifier.
\newline -- \code{postId}: \code{number | null} -- post target, or \code{null}.
\newline -- \code{commentId}: \code{number | null} -- comment target, or \code{null}.
\newline -- \textit{Returns:} \code{Promise<{ id: number } | null>} -- the report id when one exists. \\ \hline

05 & \code{getForumReportLists(}\newline\code{roleName: string, page: number,}\newline\code{limit: number, status?: string,}\newline\code{search?: string)} &
A method to authorize an administrator and return the report queue.
\newline -- \code{roleName}: \code{string} -- must equal \code{ADMIN}.
\newline -- \code{page}, \code{limit}: \code{number} -- pagination.
\newline -- \code{status}, \code{search}: \code{string | undefined} -- optional filters.
\newline -- \textit{Returns:} \code{Promise<ForumReportListResponseType>} -- reports and pagination.
\newline -- Throws \code{ForumReportForbiddenException} or \code{FailedToLoadForumReportsException}. \\ \hline

06 & \code{updateReportStatus(}\newline\code{reportId: number, status: string,}\newline\code{adminId: number, roleName: string)} &
A method to authorize an administrator and update a post report status.
\newline -- \code{reportId}: \code{number} -- report identifier.
\newline -- \code{status}: \code{string} -- target status.
\newline -- \code{adminId}: \code{number} -- administrator identifier.
\newline -- \code{roleName}: \code{string} -- must equal \code{ADMIN}.
\newline -- \textit{Returns:} \code{Promise<UpdateReportStatusResponseType>} -- the updated report.
\newline -- Throws \code{ForumReportForbiddenException}, \code{ForumReportNotFoundException}, or \code{FailedToUpdateReportStatusException}. \\ \hline
\end{ClassMethodTable}

\paragraph{ClassAttribute}
\begin{ClassAttributeTable}
01 & \code{forumReportRepository} & \code{ForumReportRepository} -- private readonly persistence dependency. \\ \hline
\end{ClassAttributeTable}

\subsection{ForumReportRepository Class}

\paragraph{ClassMethod}
\begin{ClassMethodTable}
01 & \code{ForumReportRepository(prisma: PrismaService)} &
A method to initialize the repository with the shared Prisma service.
\newline -- \code{prisma}: \code{PrismaService} -- database client.
\newline -- \textit{Returns:} \code{void} -- constructor, no return value. \\ \hline

02 & \code{findPostById(postId: number)} &
A method to check that a non-deleted post exists.
\newline -- \code{postId}: \code{number} -- post identifier.
\newline -- \textit{Returns:} \code{Promise<{ id: number } | null>} -- the id when found, otherwise \code{null}. \\ \hline

03 & \code{findCommentById(commentId: number)} &
A method to check that a non-deleted comment exists and belongs to a post.
\newline -- \code{commentId}: \code{number} -- comment identifier.
\newline -- \textit{Returns:} \code{Promise<{ id: number; postId: number } | null>} -- comment id and parent post id. \\ \hline

04 & \code{findExistingReport(}\newline\code{reporterId: number, postId: number | null,}\newline\code{commentId: number | null)} &
A method to find a duplicate report by the same user and target.
\newline -- \code{reporterId}: \code{number} -- reporter identifier.
\newline -- \code{postId}: \code{number | null} -- post target, or \code{null}.
\newline -- \code{commentId}: \code{number | null} -- comment target, or \code{null}.
\newline -- \textit{Returns:} \code{Promise<{ id: number } | null>} -- the report id when one exists. \\ \hline

05 & \code{createReport(}\newline\code{reporterId: number, postId: number | null,}\newline\code{commentId: number | null, reason: string)} &
A method to insert a report.
\newline -- \code{reporterId}: \code{number} -- reporter identifier.
\newline -- \code{postId}, \code{commentId}: \code{number | null} -- exactly one target is populated by the service.
\newline -- \code{reason}: \code{string} -- report reason.
\newline -- \textit{Returns:} \code{Promise<ForumReportType>} -- the created report. \\ \hline

06 & \code{findReportById(id: number)} &
A method to load a report by identifier.
\newline -- \code{id}: \code{number} -- report identifier.
\newline -- \textit{Returns:} \code{Promise<ForumReportType | null>} -- the report or \code{null}. \\ \hline

07 & \code{getReportList(}\newline\code{page: number, limit: number,}\newline\code{status?: string, search?: string)} &
A method to query a page of reports and the matching total in one transaction.
\newline -- \code{page}, \code{limit}: \code{number} -- pagination.
\newline -- \code{status}: \code{string | undefined} -- optional status filter.
\newline -- \code{search}: \code{string | undefined} -- optional case-insensitive reason search.
\newline -- \textit{Returns:} \code{Promise<{ reports: ForumReportListResponseType['reports']; total: number }>} -- report rows and total. \\ \hline

08 & \code{updateReportStatus(}\newline\code{id: number, status: string,}\newline\code{adminId: number | null)} &
A method to update a report status and record the reviewing administrator.
\newline -- \code{id}: \code{number} -- report identifier.
\newline -- \code{status}: \code{string} -- target status, converted to upper case.
\newline -- \code{adminId}: \code{number | null} -- reviewing administrator.
\newline -- \textit{Returns:} \code{Promise<ForumReportType>} -- the updated report. \\ \hline
\end{ClassMethodTable}

\paragraph{ClassAttribute}
\begin{ClassAttributeTable}
01 & \code{prisma} & \code{PrismaService} -- private readonly database client. \\ \hline
\end{ClassAttributeTable}

\subsubsection*{Supporting DTO classes for reports}
\addcontentsline{toc}{subsubsection}{Supporting DTO classes for reports}

\begin{SupportingTable}{4.4cm}{7.2cm}{4.2cm}
\code{CreateForumReportDto} & \code{reason: string}. & Report-creation body. \\ \hline
\code{CreateForumReportResponseDto} & \code{ForumReportType}: \code{id}, \code{reporterId}, \code{postId}, \code{commentId}, \code{reason}, \code{status}, \code{adminId}, \code{createdAt}. & Report-creation response. \\ \hline
\code{UpdateReportStatusDto} & \code{status: PENDING | REVIEWED | RESOLVED | DISMISSED}. & Report-status body. \\ \hline
\code{UpdateReportStatusResponseDto} & \code{ForumReportType}. & Report-status response. \\ \hline
\code{ForumReportListResponseDto} & \code{reports: ForumReportType[]} with reporter, post, and comment summaries, plus \code{pagination}. & Administrator report queue. \\ \hline
\end{SupportingTable}

% ========================================================
% SPRINT 2
% ========================================================
% =================================================================
% FILE: sections/kvy/4/sprint2.tex -- Sprint 2 (Conversation/Messaging)
% Class specifications for the direct-conversation implementation.
% =================================================================
\subsection{ConversationController Class}

\paragraph{ClassMethod}
\begin{ClassMethodTable}
01 & \code{ConversationController(}\newline\code{conversationService: ConversationService)} &
A method to initialize the controller with the conversation service.
\newline -- \code{conversationService}: \code{ConversationService} -- service that owns all conversation use cases.
\newline -- \textit{Returns:} \code{void} -- constructor, no return value. \\ \hline

02 & \code{getConversations(userId: number)} &
A method to return the current user's conversation list.
\newline -- \code{userId}: \code{number} -- current user identifier injected by \code{@UserActive('userId')}.
\newline -- \textit{Returns:} \code{Promise<GetConversationsResponseType>} -- conversation items with other user, last message, unread count, and pin time.
\newline -- Route: \code{GET /conversations}; authenticated endpoint. \\ \hline

03 & \code{createConversation(}\newline\code{userId: number, body: CreateConversationBodyType)} &
A method to create or reopen a conversation with another user.
\newline -- \code{userId}: \code{number} -- current user identifier.
\newline -- \code{body.participantId}: \code{number} -- the other participant, validated by \code{CreateConversationBodyDto}.
\newline -- \textit{Returns:} \code{Promise<CreateConversationResponseType>} -- id, participant1, participant2, and createdAt.
\newline -- Route: \code{POST /conversations}; authenticated endpoint. \\ \hline

04 & \code{getMessages(}\newline\code{userId: number, id: number)} &
A method to return the message history of one conversation.
\newline -- \code{userId}: \code{number} -- current user identifier.
\newline -- \code{id}: \code{number} -- conversation identifier parsed by \code{ParseIntPipe}.
\newline -- \textit{Returns:} \code{Promise<DirectMessageType[]>} -- non-deleted messages in ascending creation order.
\newline -- Route: \code{GET /conversations/:id/messages}; authenticated endpoint. \\ \hline

05 & \code{sendMessage(}\newline\code{userId: number, id: number,}\newline\code{body: SendMessageBodyType)} &
A method to send a message through the REST fallback.
\newline -- \code{userId}: \code{number} -- sender identifier.
\newline -- \code{id}: \code{number} -- conversation identifier.
\newline -- \code{body}: \code{SendMessageBodyType} -- message text and optional attachment, validated by \code{SendMessageBodyDto}.
\newline -- \textit{Returns:} \code{Promise<DirectMessageType>} -- the saved message.
\newline -- Route: \code{POST /conversations/:id/messages}; authenticated endpoint. \\ \hline

06 & \code{uploadFile(}\newline\code{userId: number, id: number,}\newline\code{file?: Express.Multer.File)} &
A method to upload a message attachment for a conversation participant.
\newline -- \code{userId}: \code{number} -- current user identifier.
\newline -- \code{id}: \code{number} -- conversation identifier.
\newline -- \code{file}: \code{Express.Multer.File | undefined} -- multipart file read into memory by \code{FileInterceptor}.
\newline -- \textit{Returns:} \code{Promise<MessageAttachmentType>} -- Cloudinary URL, original name, byte size, and MIME type.
\newline -- Route: \code{POST /conversations/:id/files}; authenticated endpoint. \\ \hline

07 & \code{deleteMessage(}\newline\code{userId: number, id: number,}\newline\code{messageId: number)} &
A method to soft-delete a message sent by the current user.
\newline -- \code{userId}: \code{number} -- current user identifier.
\newline -- \code{id}: \code{number} -- conversation identifier.
\newline -- \code{messageId}: \code{number} -- message identifier.
\newline -- \textit{Returns:} \code{Promise<DeleteMessageResponseType>} -- the deleted message id.
\newline -- Route: \code{DELETE /conversations/:id/messages/:messageId}; authenticated endpoint. \\ \hline

08 & \code{markAsRead(}\newline\code{userId: number, id: number)} &
A method to mark all messages from the other participant as read.
\newline -- \code{userId}: \code{number} -- current user identifier.
\newline -- \code{id}: \code{number} -- conversation identifier.
\newline -- \textit{Returns:} \code{Promise<{ count: number }>} -- number of updated messages. This route has no response DTO decorator.
\newline -- Route: \code{PATCH /conversations/:id/read}; authenticated endpoint. \\ \hline

09 & \code{hideConversation(}\newline\code{userId: number, id: number)} &
A method to hide a conversation only for the current user.
\newline -- \code{userId}: \code{number} -- current user identifier.
\newline -- \code{id}: \code{number} -- conversation identifier.
\newline -- \textit{Returns:} \code{Promise<HideConversationResponseType>} -- id and hidden flag.
\newline -- Route: \code{PATCH /conversations/:id/hide}; authenticated endpoint. \\ \hline

10 & \code{pinConversation(}\newline\code{userId: number, id: number,}\newline\code{body: PinConversationBodyType)} &
A method to pin or unpin a conversation for the current user.
\newline -- \code{userId}: \code{number} -- current user identifier.
\newline -- \code{id}: \code{number} -- conversation identifier.
\newline -- \code{body.pinned}: \code{boolean} -- target pin state, validated by \code{PinConversationBodyDto}.
\newline -- \textit{Returns:} \code{Promise<PinConversationResponseType>} -- id and the new \code{pinnedAt} value.
\newline -- Route: \code{PATCH /conversations/:id/pin}; authenticated endpoint. \\ \hline
\end{ClassMethodTable}

\paragraph{ClassAttribute}
\begin{ClassAttributeTable}
01 & \code{conversationService} & \code{ConversationService} -- private readonly service dependency. \\ \hline
\end{ClassAttributeTable}

\subsection{ConversationService Class}

\paragraph{ClassMethod}
\begin{ClassMethodTable}
01 & \code{ConversationService(}\newline\code{conversationRepository: ConversationRepository)} &
A method to initialize the service with the conversation repository.
\newline -- \code{conversationRepository}: \code{ConversationRepository} -- persistence dependency.
\newline -- \textit{Returns:} \code{void} -- constructor, no return value. \\ \hline

02 & \code{buildConversationItem(}\newline\code{conversation: RawConversationItem,}\newline\code{userId: number)} &
A private method to map a repository row to the public conversation item.
\newline -- \code{conversation}: \code{RawConversationItem} -- row with both participants, last message, and unread count.
\newline -- \code{userId}: \code{number} -- current user, used to select the other participant and pin time.
\newline -- \textit{Returns:} \code{ConversationListItemType} -- normalized item with \code{otherUser}, \code{lastMessage}, \code{unreadCount}, and \code{pinnedAt}. \\ \hline

03 & \code{getConversations(userId: number)} &
A method to load and sort all visible conversations of a user.
\newline -- \code{userId}: \code{number} -- current user identifier.
\newline -- \textit{Returns:} \code{Promise<GetConversationsResponseType>} -- pinned conversations first, then by last-message time or creation time.
\newline -- Any failure is mapped to \code{FailedToLoadConversationsException}. \\ \hline

04 & \code{getConversationItem(}\newline\code{conversationId: number, userId: number)} &
A method to rebuild one conversation item for a realtime update.
\newline -- \code{conversationId}: \code{number} -- conversation identifier.
\newline -- \code{userId}: \code{number} -- recipient of the realtime item.
\newline -- \textit{Returns:} \code{Promise<ConversationListItemType | null>} -- the item, or \code{null} when the conversation is not visible to that user. \\ \hline

05 & \code{getMessages(}\newline\code{userId: number, conversationId: number)} &
A method to authorize a participant and return the message history.
\newline -- \code{userId}: \code{number} -- current user identifier.
\newline -- \code{conversationId}: \code{number} -- conversation identifier.
\newline -- \textit{Returns:} \code{Promise<DirectMessageType[]>} -- non-deleted messages.
\newline -- Throws \code{ConversationNotFoundException} when the user is not a participant; other failures become \code{FailedToLoadConversationsException}. \\ \hline

06 & \code{createConversation(}\newline\code{userId: number, participantId: number)} &
A method to create a conversation or return the existing pair.
\newline -- \code{userId}: \code{number} -- current user identifier.
\newline -- \code{participantId}: \code{number} -- the other user.
\newline -- \textit{Returns:} \code{Promise<CreateConversationResponseType>} -- id, participant1, participant2, and createdAt.
\newline -- Rejects self-conversation with \code{CannotStartConversationWithSelfException}; an existing hidden conversation is unhidden. \\ \hline

07 & \code{sendMessage(}\newline\code{userId: number, conversationId: number,}\newline\code{message: string,}\newline\code{attachment?: MessageAttachmentType | null)} &
A method to save a message and unhide the conversation for both participants.
\newline -- \code{userId}: \code{number} -- sender identifier.
\newline -- \code{conversationId}: \code{number} -- conversation identifier.
\newline -- \code{message}: \code{string} -- text; may be empty when an attachment exists.
\newline -- \code{attachment}: \code{MessageAttachmentType | null | undefined} -- optional file metadata.
\newline -- \textit{Returns:} \code{Promise<DirectMessageType>} -- the saved message.
\newline -- Throws \code{ConversationNotFoundException} or \code{FailedToSendMessageException}. \\ \hline

08 & \code{uploadFile(}\newline\code{userId: number, conversationId: number,}\newline\code{file?: Express.Multer.File)} &
A method to validate participation and upload a file to Cloudinary.
\newline -- \code{userId}: \code{number} -- current user identifier.
\newline -- \code{conversationId}: \code{number} -- conversation identifier.
\newline -- \code{file}: \code{Express.Multer.File | undefined} -- in-memory multipart file.
\newline -- \textit{Returns:} \code{Promise<MessageAttachmentType>} -- secure URL, original name, byte size, and MIME type.
\newline -- Throws \code{FailedToUploadFileException}, \code{ConversationNotFoundException}, or passes through a Nest \code{HttpException}. \\ \hline

09 & \code{deleteMessage(}\newline\code{userId: number, conversationId: number,}\newline\code{messageId: number)} &
A method to authorize a participant and soft-delete the sender's message.
\newline -- \code{userId}: \code{number} -- current user identifier; must be the sender.
\newline -- \code{conversationId}: \code{number} -- conversation identifier.
\newline -- \code{messageId}: \code{number} -- message identifier.
\newline -- \textit{Returns:} \code{Promise<{ id: number }>} -- the deleted message id.
\newline -- Throws \code{ConversationNotFoundException}, \code{MessageNotFoundException}, or \code{FailedToDeleteMessageException}. \\ \hline

10 & \code{markAsRead(}\newline\code{userId: number, conversationId: number)} &
A method to mark messages from the other participant as read.
\newline -- \code{userId}: \code{number} -- current user identifier.
\newline -- \code{conversationId}: \code{number} -- conversation identifier.
\newline -- \textit{Returns:} \code{Promise<{ count: number }>} -- number of updated messages.
\newline -- Throws \code{ConversationNotFoundException} or \code{FailedToMarkAsReadException}. \\ \hline

11 & \code{hideConversation(}\newline\code{userId: number, conversationId: number)} &
A method to set the current participant's hidden timestamp.
\newline -- \code{userId}: \code{number} -- current user identifier.
\newline -- \code{conversationId}: \code{number} -- conversation identifier.
\newline -- \textit{Returns:} \code{Promise<HideConversationResponseType>} -- id and hidden flag.
\newline -- Throws \code{ConversationNotFoundException} or \code{FailedToHideConversationException}. \\ \hline

12 & \code{pinConversation(}\newline\code{userId: number, conversationId: number,}\newline\code{pinned: boolean)} &
A method to set or clear the current participant's pin timestamp.
\newline -- \code{userId}: \code{number} -- current user identifier.
\newline -- \code{conversationId}: \code{number} -- conversation identifier.
\newline -- \code{pinned}: \code{boolean} -- true to pin, false to unpin.
\newline -- \textit{Returns:} \code{Promise<PinConversationResponseType>} -- id and \code{pinnedAt}.
\newline -- Throws \code{ConversationNotFoundException} or \code{FailedToPinConversationException}. \\ \hline

13 & \code{ensureParticipant(}\newline\code{conversationId: number, userId: number)} &
A method used by the gateway to assert that a socket user belongs to a conversation.
\newline -- \code{conversationId}: \code{number} -- conversation identifier.
\newline -- \code{userId}: \code{number} -- socket user identifier.
\newline -- \textit{Returns:} the non-deleted \code{Conversation} row.
\newline -- Throws \code{ConversationAccessDeniedException} when the user is not a participant. \\ \hline
\end{ClassMethodTable}

\paragraph{ClassAttribute}
\begin{ClassAttributeTable}
01 & \code{conversationRepository} & \code{ConversationRepository} -- private readonly persistence dependency. \\ \hline
02 & \code{logger} & \code{Logger} -- private readonly NestJS logger; records attachment-upload failures. \\ \hline
\end{ClassAttributeTable}

\subsection{ConversationRepository Class}

\paragraph{ClassMethod}
\begin{ClassMethodTable}
01 & \code{ConversationRepository(prisma: PrismaService)} &
A method to initialize the repository with the shared Prisma service.
\newline -- \code{prisma}: \code{PrismaService} -- database client.
\newline -- \textit{Returns:} \code{void} -- constructor, no return value. \\ \hline

02 & \code{findConversationBetween(}\newline\code{userId1: number, userId2: number)} &
A method to find the non-deleted conversation for a user pair in either participant order.
\newline -- \code{userId1}, \code{userId2}: \code{number} -- the two users.
\newline -- \textit{Returns:} \code{Promise<Conversation | null>} -- the conversation or \code{null}. \\ \hline

03 & \code{getConversationsByUserId(userId: number)} &
A method to query all conversations visible to a user with previews and unread counts.
\newline -- \code{userId}: \code{number} -- current user identifier.
\newline -- \textit{Returns:} conversation rows with both participant profiles, the latest non-deleted message, and a filtered message count.
\newline -- A conversation is included only when the current participant's hidden timestamp is \code{null}. \\ \hline

04 & \code{findConversationByIdForUser(}\newline\code{conversationId: number, userId: number)} &
A method to load a non-deleted conversation only when the user is a participant.
\newline -- \code{conversationId}: \code{number} -- conversation identifier.
\newline -- \code{userId}: \code{number} -- current user identifier.
\newline -- \textit{Returns:} \code{Promise<Conversation | null>} -- the conversation or \code{null}. \\ \hline

05 & \code{findConversationById(conversationId: number)} &
A method to load a non-deleted conversation without a participant check.
\newline -- \code{conversationId}: \code{number} -- conversation identifier.
\newline -- \textit{Returns:} \code{Promise<Conversation | null>} -- the conversation or \code{null}. \\ \hline

06 & \code{hideConversationForUser(}\newline\code{conversationId: number, userId: number)} &
A method to set the hidden timestamp of the participant who requested the hide.
\newline -- \code{conversationId}: \code{number} -- conversation identifier.
\newline -- \code{userId}: \code{number} -- participant identifier.
\newline -- \textit{Returns:} \code{Promise<boolean>} -- true when the row was updated, false when the conversation or participant does not match. \\ \hline

07 & \code{unhideConversation(conversationId: number)} &
A method to clear both participants' hidden timestamps after new activity.
\newline -- \code{conversationId}: \code{number} -- conversation identifier.
\newline -- \textit{Returns:} the updated \code{Conversation} row. \\ \hline

08 & \code{setPinnedForUser(}\newline\code{conversationId: number, userId: number,}\newline\code{pinned: boolean)} &
A method to set or clear the participant's pin timestamp.
\newline -- \code{conversationId}: \code{number} -- conversation identifier.
\newline -- \code{userId}: \code{number} -- participant identifier.
\newline -- \code{pinned}: \code{boolean} -- true sets the current time, false clears it.
\newline -- \textit{Returns:} \code{Promise<Date | null>} -- the new timestamp, or \code{null} when unpinned or not a participant. \\ \hline

09 & \code{getMessages(conversationId: number)} &
A method to query all non-deleted messages of a conversation in ascending time order.
\newline -- \code{conversationId}: \code{number} -- conversation identifier.
\newline -- \textit{Returns:} \code{Promise<DirectMessage[]>} -- Prisma message rows. \\ \hline

10 & \code{createConversation(}\newline\code{participant1: number, participant2: number)} &
A method to insert a conversation pair.
\newline -- \code{participant1}: \code{number} -- first participant (the creator in the service flow).
\newline -- \code{participant2}: \code{number} -- second participant.
\newline -- \textit{Returns:} \code{Promise<Conversation>} -- the inserted row. \\ \hline

11 & \code{createMessage(}\newline\code{conversationId: number, senderId: number,}\newline\code{message: string,}\newline\code{attachment?: MessageAttachmentType)} &
A method to insert a message and copy optional attachment metadata.
\newline -- \code{conversationId}: \code{number} -- conversation identifier.
\newline -- \code{senderId}: \code{number} -- sender identifier.
\newline -- \code{message}: \code{string} -- message text.
\newline -- \code{attachment}: \code{MessageAttachmentType | undefined} -- optional file metadata.
\newline -- \textit{Returns:} \code{Promise<DirectMessage>} -- the inserted message. \\ \hline

12 & \code{softDeleteMessage(}\newline\code{conversationId: number, messageId: number,}\newline\code{senderId: number)} &
A method to set \code{deletedAt} only when the message belongs to the conversation and sender.
\newline -- \code{conversationId}: \code{number} -- conversation identifier.
\newline -- \code{messageId}: \code{number} -- message identifier.
\newline -- \code{senderId}: \code{number} -- sender identifier.
\newline -- \textit{Returns:} \code{Promise<boolean>} -- true when one non-deleted row was updated. \\ \hline

13 & \code{markMessagesAsRead(}\newline\code{conversationId: number, userId: number)} &
A method to mark messages from the other participant as read.
\newline -- \code{conversationId}: \code{number} -- conversation identifier.
\newline -- \code{userId}: \code{number} -- current user; messages sent by this user are excluded.
\newline -- \textit{Returns:} the Prisma batch payload with the updated row count. \\ \hline

14 & \code{countUnread(}\newline\code{conversationId: number, userId: number)} &
A method to count unread messages from the other participant.
\newline -- \code{conversationId}: \code{number} -- conversation identifier.
\newline -- \code{userId}: \code{number} -- current user.
\newline -- \textit{Returns:} \code{Promise<number>} -- count of non-deleted, unread messages not sent by the user. \\ \hline
\end{ClassMethodTable}

\paragraph{ClassAttribute}
\begin{ClassAttributeTable}
01 & \code{prisma} & \code{PrismaService} -- private readonly database client. \\ \hline
\end{ClassAttributeTable}

\subsection{ConversationsGateway Class}

\paragraph{ClassMethod}
\begin{ClassMethodTable}
01 & \code{ConversationsGateway(}\newline\code{conversationService: ConversationService,}\newline\code{tokenService: TokenService)} &
A method to initialize the gateway dependencies.
\newline -- \code{conversationService}: \code{ConversationService} -- shared use-case service.
\newline -- \code{tokenService}: \code{TokenService} -- verifies the access token during the socket handshake.
\newline -- \textit{Returns:} \code{void} -- constructor, no return value. \\ \hline

02 & \code{handleConnection(client: Socket)} &
A method to authenticate a new socket and join its user room.
\newline -- \code{client}: \code{Socket} -- connecting Socket.IO client.
\newline -- Reads \code{client.handshake.auth.token}, falling back to the bearer token in the authorization header; verifies it with \code{TokenService.verifyAccessToken}.
\newline -- \textit{Returns:} \code{Promise<void>} -- stores \code{userId} in \code{client.data} and joins \code{user:<userId>}; disconnects on failure. \\ \hline

03 & \code{handleDisconnect(client: Socket)} &
A method to log a disconnected socket.
\newline -- \code{client}: \code{Socket} -- disconnected client.
\newline -- \textit{Returns:} \code{void} -- no return value. \\ \hline

04 & \code{handleSendMessage(}\newline\code{client: Socket, payload: unknown)} &
A method to handle the \code{message:send} event.
\newline -- \code{client}: \code{Socket} -- sender socket.
\newline -- \code{payload}: \code{unknown} -- validated by \code{SendMessagePayloadSchema}.
\newline -- Ensures participation, saves the message, emits \code{message:new} to both user rooms, and refreshes both conversation lists.
\newline -- \textit{Returns:} \code{Promise<{ success: boolean; message?: DirectMessageType }>} -- success flag and saved message; emits \code{message:error} on failure. \\ \hline

05 & \code{handleMarkAsRead(}\newline\code{client: Socket, payload: unknown)} &
A method to handle the \code{message:read} event.
\newline -- \code{client}: \code{Socket} -- reader socket.
\newline -- \code{payload}: \code{unknown} -- validated by \code{MarkAsReadPayloadSchema}.
\newline -- Calls \code{markAsRead} and emits \code{conversation:updated} to both participants.
\newline -- \textit{Returns:} \code{Promise<{ success: boolean }>} -- success flag. \\ \hline

06 & \code{handleDeleteMessage(}\newline\code{client: Socket, payload: unknown)} &
A method to handle the \code{message:delete} event.
\newline -- \code{client}: \code{Socket} -- sender socket.
\newline -- \code{payload}: \code{unknown} -- validated by \code{DeleteMessagePayloadSchema}.
\newline -- Calls \code{deleteMessage}, emits \code{message:deleted} to both user rooms, and refreshes both conversation lists.
\newline -- \textit{Returns:} \code{Promise<{ success: boolean }>} -- success flag; emits \code{message:error} on failure. \\ \hline

07 & \code{handleHideConversation(}\newline\code{client: Socket, payload: unknown)} &
A method to handle the \code{conversation:hide} event.
\newline -- \code{client}: \code{Socket} -- current user's socket.
\newline -- \code{payload}: \code{unknown} -- validated by \code{HideConversationPayloadSchema}.
\newline -- Calls \code{hideConversation} and emits \code{conversation:hidden} only to the current user's room.
\newline -- \textit{Returns:} \code{Promise<{ success: boolean }>} -- success flag. \\ \hline

08 & \code{handleTyping(}\newline\code{client: Socket, payload: unknown)} &
A method to handle the \code{typing} event.
\newline -- \code{client}: \code{Socket} -- sender socket.
\newline -- \code{payload}: \code{unknown} -- validated by \code{TypingPayloadSchema}.
\newline -- Ensures participation and forwards \code{typing} with \code{conversationId} and \code{isTyping} only to the other user's room.
\newline -- \textit{Returns:} \code{Promise<{ success: boolean }>} -- success flag. \\ \hline

09 & \code{emitConversationUpdated(}\newline\code{conversationId: number, userId: number)} &
A private method to push the latest conversation item to one user.
\newline -- \code{conversationId}: \code{number} -- conversation identifier.
\newline -- \code{userId}: \code{number} -- target user.
\newline -- \textit{Returns:} \code{Promise<void>} -- emits \code{conversation:updated} when \code{getConversationItem} returns a visible item. \\ \hline
\end{ClassMethodTable}

\paragraph{ClassAttribute}
\begin{ClassAttributeTable}
01 & \code{server} & \code{Server} -- Socket.IO server injected by \code{@WebSocketServer()}; used to address user rooms. \\ \hline
02 & \code{logger} & \code{Logger} -- private readonly NestJS logger for connection and event failures. \\ \hline
03 & \code{conversationService} & \code{ConversationService} -- private readonly service dependency. \\ \hline
04 & \code{tokenService} & \code{TokenService} -- private readonly access-token verifier. \\ \hline
\end{ClassAttributeTable}

% ========================================================
% SPRINT 3
% ========================================================
% =================================================================
% FILE: sections/kvy/4/sprint3.tex -- Sprint 3 (Administration)
% Class specifications for forum and identity-verification administration.
% =================================================================.

\subsection{ForumAdminController Class}

\paragraph{ClassMethod}
\begin{ClassMethodTable}
01 & \code{ForumAdminController(}\newline\code{adminService: ForumAdminService)} &
A method to initialize the controller with the administration service.
\newline -- \code{adminService}: \code{ForumAdminService} -- service dependency.
\newline -- \textit{Returns:} \code{void} -- constructor, no return value. \\ \hline

02 & \code{getAdminStats(roleName: string)} &
A method to return the forum dashboard counters.
\newline -- \code{roleName}: \code{string} -- current role name from \code{@UserActive('roleName')}.
\newline -- \textit{Returns:} \code{Promise<ForumAdminStatsType>} -- category, post, comment, report, pending-report, and user totals plus the five newest posts.
\newline -- Route: \code{GET /forums/admin/stats}; administrator endpoint. \\ \hline

03 & \code{getAdminCategoryLists(}\newline\code{roleName: string, page: number,}\newline\code{limit: number, search: string,}\newline\code{deleted: string,}\newline\code{sortBy: 'id' | 'name' | 'createdAt',}\newline\code{sortOrder: 'asc' | 'desc')} &
A method to list categories with search, deleted-state filtering, and sorting.
\newline -- \code{roleName}: \code{string} -- current role name.
\newline -- \code{page}, \code{limit}: \code{number} -- coerced with \code{Number()} and defaulted to 1 and 10.
\newline -- \code{search}: \code{string} -- optional case-insensitive name search.
\newline -- \code{deleted}: \code{string} -- optional \code{'true'} or \code{'false'} filter.
\newline -- \code{sortBy}, \code{sortOrder}: \code{'id' | 'name' | 'createdAt'} and \code{'asc' | 'desc'} -- optional sort controls.
\newline -- \textit{Returns:} \code{Promise<ForumAdminCategoryListResponseType>} -- categories and pagination.
\newline -- Route: \code{GET /forums/admin/categories}; administrator endpoint. \\ \hline

04 & \code{getAdminCategoryById(}\newline\code{roleName: string, id: number)} &
A method to return one active category.
\newline -- \code{roleName}: \code{string} -- current role name.
\newline -- \code{id}: \code{number} -- category identifier.
\newline -- \textit{Returns:} \code{Promise<ForumCategoryDetailResponseType>} -- category detail with public post count.
\newline -- Route: \code{GET /forums/admin/categories/:id}; administrator endpoint. \\ \hline

05 & \code{createAdminCategory(}\newline\code{roleName: string,}\newline\code{body: CreateForumCategoryBodyDto)} &
A method to create a category.
\newline -- \code{roleName}: \code{string} -- current role name.
\newline -- \code{body}: \code{CreateForumCategoryBodyDto} -- name and optional description.
\newline -- \textit{Returns:} \code{Promise<ForumCategoryType>} -- the created category.
\newline -- Route: \code{POST /forums/admin/categories}; administrator endpoint. \\ \hline

06 & \code{updateAdminCategory(}\newline\code{roleName: string, id: number,}\newline\code{body: UpdateForumCategoryBodyDto)} &
A method to rename or re-describe a category.
\newline -- \code{roleName}: \code{string} -- current role name.
\newline -- \code{id}: \code{number} -- category identifier.
\newline -- \code{body}: \code{UpdateForumCategoryBodyDto} -- optional name and optional description.
\newline -- \textit{Returns:} \code{Promise<ForumCategoryType>} -- the updated category.
\newline -- Route: \code{PATCH /forums/admin/categories/:id}; administrator endpoint. \\ \hline

07 & \code{deleteAdminCategory(}\newline\code{roleName: string, id: number)} &
A method to soft-delete a category and detach its posts.
\newline -- \code{roleName}: \code{string} -- current role name.
\newline -- \code{id}: \code{number} -- category identifier.
\newline -- \textit{Returns:} \code{Promise<{ message: string }>} -- deletion confirmation.
\newline -- Route: \code{DELETE /forums/admin/categories/:id}; administrator endpoint. \\ \hline

08 & \code{restoreAdminCategory(}\newline\code{roleName: string, id: number)} &
A method to restore a soft-deleted category.
\newline -- \code{roleName}: \code{string} -- current role name.
\newline -- \code{id}: \code{number} -- category identifier.
\newline -- \textit{Returns:} \code{Promise<ForumCategoryType>} -- the restored category.
\newline -- Route: \code{PATCH /forums/admin/categories/:id/restore}; administrator endpoint. \\ \hline

09 & \code{getAdminCommentLists(}\newline\code{roleName: string, page: number,}\newline\code{limit: number, search: string)} &
A method to list active comments for administration.
\newline -- \code{roleName}: \code{string} -- current role name.
\newline -- \code{page}, \code{limit}: \code{number} -- pagination, defaulted to 1 and 10.
\newline -- \code{search}: \code{string} -- optional case-insensitive content search.
\newline -- \textit{Returns:} \code{Promise<ForumAdminCommentListResponseType>} -- comments with author and parent-post summaries, plus pagination.
\newline -- Route: \code{GET /forums/admin/comments}; administrator endpoint. \\ \hline

10 & \code{getPendingPosts(}\newline\code{roleName: string, page: number,}\newline\code{limit: number)} &
A method to list posts waiting for moderation.
\newline -- \code{roleName}: \code{string} -- current role name.
\newline -- \code{page}, \code{limit}: \code{number} -- pagination, defaulted to 1 and 10.
\newline -- \textit{Returns:} \code{Promise<PendingForumPostListResponseType>} -- non-deleted \code{PENDING} posts, oldest first, with pagination.
\newline -- Route: \code{GET /forums/admin/pending-posts}; administrator endpoint. \\ \hline

11 & \code{getTrashPosts(}\newline\code{roleName: string, page: number,}\newline\code{limit: number)} &
A method to list deleted or rejected posts.
\newline -- \code{roleName}: \code{string} -- current role name.
\newline -- \code{page}, \code{limit}: \code{number} -- pagination, defaulted to 1 and 10.
\newline -- \textit{Returns:} \code{Promise<ForumTrashPostListResponseType>} -- posts with \code{deletedAt} set or moderation status \code{REJECTED}.
\newline -- Route: \code{GET /forums/admin/trash/posts}; administrator endpoint. \\ \hline

12 & \code{getTrashComments(}\newline\code{roleName: string, page: number,}\newline\code{limit: number)} &
A method to list deleted comments.
\newline -- \code{roleName}: \code{string} -- current role name.
\newline -- \code{page}, \code{limit}: \code{number} -- pagination, defaulted to 1 and 10.
\newline -- \textit{Returns:} \code{Promise<ForumTrashCommentListResponseType>} -- soft-deleted comments with author and parent-post summaries.
\newline -- Route: \code{GET /forums/admin/trash/comments}; administrator endpoint. \\ \hline

13 & \code{restoreTrashPost(}\newline\code{roleName: string, id: number)} &
A method to restore a post from the trash.
\newline -- \code{roleName}: \code{string} -- current role name.
\newline -- \code{id}: \code{number} -- post identifier.
\newline -- \textit{Returns:} \code{Promise<ForumPostType>} -- the restored post.
\newline -- Route: \code{PATCH /forums/admin/trash/posts/:id/restore}; administrator endpoint. \\ \hline

14 & \code{restoreTrashComment(}\newline\code{roleName: string, id: number)} &
A method to restore a comment from the trash.
\newline -- \code{roleName}: \code{string} -- current role name.
\newline -- \code{id}: \code{number} -- comment identifier.
\newline -- \textit{Returns:} \code{Promise<ForumAdminCommentType>} -- the restored comment.
\newline -- Route: \code{PATCH /forums/admin/trash/comments/:id/restore}; administrator endpoint. \\ \hline

15 & \code{approvePendingPost(}\newline\code{roleName: string, userId: number,}\newline\code{id: number)} &
A method to approve a pending post.
\newline -- \code{roleName}: \code{string} -- current role name.
\newline -- \code{userId}: \code{number} -- reviewing administrator recorded in \code{reviewedById}.
\newline -- \code{id}: \code{number} -- post identifier.
\newline -- \textit{Returns:} \code{Promise<ForumPostType>} -- the approved post.
\newline -- Route: \code{PATCH /forums/admin/posts/:id/approve}; administrator endpoint. \\ \hline

16 & \code{rejectPendingPost(}\newline\code{roleName: string, userId: number,}\newline\code{id: number)} &
A method to reject a pending post and move it to the trash.
\newline -- \code{roleName}: \code{string} -- current role name.
\newline -- \code{userId}: \code{number} -- reviewing administrator.
\newline -- \code{id}: \code{number} -- post identifier.
\newline -- \textit{Returns:} \code{Promise<ForumPostType>} -- the rejected post, now soft-deleted.
\newline -- Route: \code{PATCH /forums/admin/posts/:id/reject}; administrator endpoint. \\ \hline
\end{ClassMethodTable}

\paragraph{ClassAttribute}
\begin{ClassAttributeTable}
01 & \code{adminService} & \code{ForumAdminService} -- private readonly service dependency. \\ \hline
\end{ClassAttributeTable}

\subsection{ForumAdminService Class}

\paragraph{ClassMethod}
\begin{ClassMethodTable}
01 & \code{ForumAdminService(}\newline\code{adminRepository: ForumAdminRepository)} &
A method to initialize the service with the administration repository.
\newline -- \code{adminRepository}: \code{ForumAdminRepository} -- persistence dependency.
\newline -- \textit{Returns:} \code{void} -- constructor, no return value. \\ \hline

02 & \code{getAdminStats(roleName: string)} &
A method to authorize an administrator and load dashboard counters.
\newline -- \code{roleName}: \code{string} -- must equal \code{RoleName.ADMIN}.
\newline -- \textit{Returns:} \code{Promise<ForumAdminStatsType>} -- counters and the five newest posts.
\newline -- Throws \code{ForumReportForbiddenException} for a non-admin or a repository failure. \\ \hline

03 & \code{getAdminCommentLists(}\newline\code{roleName: string, page: number,}\newline\code{limit: number, search?: string)} &
A method to authorize an administrator and return a page of active comments.
\newline -- \code{roleName}: \code{string} -- must equal \code{RoleName.ADMIN}.
\newline -- \code{page}, \code{limit}: \code{number} -- pagination.
\newline -- \code{search}: \code{string | undefined} -- optional content filter.
\newline -- \textit{Returns:} \code{Promise<ForumAdminCommentListResponseType>} -- comments and pagination with \code{totalPages}. \\ \hline

04 & \code{getAdminCategoryLists(}\newline\code{roleName: string, page: number,}\newline\code{limit: number, search?: string,}\newline\code{sortBy?: 'id' | 'name' | 'createdAt',}\newline\code{sortOrder?: 'asc' | 'desc',}\newline\code{deleted?: string)} &
A method to authorize an administrator and return a page of categories.
\newline -- \code{roleName}: \code{string} -- must equal \code{RoleName.ADMIN}.
\newline -- \code{page}, \code{limit}: \code{number} -- pagination.
\newline -- \code{search}: \code{string | undefined} -- optional name filter.
\newline -- \code{sortBy}, \code{sortOrder}: optional sort controls; repository defaults are \code{id} and \code{desc}.
\newline -- \code{deleted}: \code{string | undefined} -- \code{'true'} lists only deleted, \code{'false'} only active.
\newline -- \textit{Returns:} \code{Promise<ForumAdminCategoryListResponseType>} -- categories and pagination. \\ \hline

05 & \code{getAdminCategoryById(}\newline\code{roleName: string, categoryId: number)} &
A method to authorize an administrator and load one active category.
\newline -- \code{roleName}: \code{string} -- must equal \code{RoleName.ADMIN}.
\newline -- \code{categoryId}: \code{number} -- category identifier.
\newline -- \textit{Returns:} \code{Promise<ForumCategoryDetailResponseType>} -- category detail.
\newline -- Throws \code{ForumCategoryNotFoundException} when absent, or \code{ForumReportForbiddenException} for other failures. \\ \hline

06 & \code{createAdminCategory(}\newline\code{roleName: string,}\newline\code{body: CreateForumCategoryBodyType)} &
A method to authorize an administrator and create a category.
\newline -- \code{roleName}: \code{string} -- must equal \code{RoleName.ADMIN}.
\newline -- \code{body}: \code{CreateForumCategoryBodyType} -- name and optional description.
\newline -- \textit{Returns:} \code{Promise<ForumCategoryType>} -- the created category.
\newline -- Unknown repository errors are rethrown unchanged. \\ \hline

07 & \code{updateAdminCategory(}\newline\code{roleName: string, categoryId: number,}\newline\code{body: UpdateForumCategoryBodyType)} &
A method to authorize an administrator and update a category.
\newline -- \code{roleName}: \code{string} -- must equal \code{RoleName.ADMIN}.
\newline -- \code{categoryId}: \code{number} -- category identifier.
\newline -- \code{body}: \code{UpdateForumCategoryBodyType} -- optional name and description.
\newline -- \textit{Returns:} \code{Promise<ForumCategoryType>} -- the updated category.
\newline -- Unknown repository errors are rethrown unchanged. \\ \hline

08 & \code{deleteAdminCategory(}\newline\code{roleName: string, categoryId: number)} &
A method to authorize an administrator and soft-delete a category.
\newline -- \code{roleName}: \code{string} -- must equal \code{RoleName.ADMIN}.
\newline -- \code{categoryId}: \code{number} -- category identifier.
\newline -- \textit{Returns:} \code{Promise<{ message: string }>} -- deletion confirmation.
\newline -- Throws \code{FailedToDeleteForumCategoryException} for an unexpected failure. \\ \hline

09 & \code{restoreAdminCategory(}\newline\code{roleName: string, categoryId: number)} &
A method to authorize an administrator and restore a deleted category.
\newline -- \code{roleName}: \code{string} -- must equal \code{RoleName.ADMIN}.
\newline -- \code{categoryId}: \code{number} -- category identifier.
\newline -- \textit{Returns:} \code{Promise<ForumCategoryType>} -- the restored category.
\newline -- Throws \code{FailedToRestoreForumCategoryException} for an unexpected failure. \\ \hline

10 & \code{getPendingPosts(}\newline\code{roleName: string, page: number,}\newline\code{limit: number)} &
A method to authorize an administrator and load the moderation queue.
\newline -- \code{roleName}: \code{string} -- must equal \code{RoleName.ADMIN}.
\newline -- \code{page}, \code{limit}: \code{number} -- pagination.
\newline -- \textit{Returns:} \code{Promise<PendingForumPostListResponseType>} -- pending posts and pagination.
\newline -- Throws \code{ForumReportForbiddenException} for a non-admin or a repository failure. \\ \hline

11 & \code{reviewPendingPost(}\newline\code{roleName: string, postId: number,}\newline\code{status: 'APPROVED' | 'REJECTED',}\newline\code{adminId: number)} &
A method to authorize an administrator and apply a moderation decision.
\newline -- \code{roleName}: \code{string} -- must equal \code{RoleName.ADMIN}.
\newline -- \code{postId}: \code{number} -- post identifier.
\newline -- \code{status}: \code{'APPROVED' | 'REJECTED'} -- decision passed by the approve/reject routes.
\newline -- \code{adminId}: \code{number} -- reviewing administrator.
\newline -- \textit{Returns:} \code{Promise<ForumPostType>} -- the reviewed post.
\newline -- A Prisma \code{P2025} error is mapped to \code{ForumPostNotFoundException}; other Prisma errors become \code{FailedToReviewForumPostException}. \\ \hline

12 & \code{getTrashPosts(}\newline\code{roleName: string, page: number,}\newline\code{limit: number)} &
A method to authorize an administrator and load the post trash.
\newline -- \code{roleName}: \code{string} -- must equal \code{RoleName.ADMIN}.
\newline -- \code{page}, \code{limit}: \code{number} -- pagination.
\newline -- \textit{Returns:} \code{Promise<ForumTrashPostListResponseType>} -- deleted or rejected posts and pagination.
\newline -- Throws \code{ForumReportForbiddenException} for a non-admin or a repository failure. \\ \hline

13 & \code{getTrashComments(}\newline\code{roleName: string, page: number,}\newline\code{limit: number)} &
A method to authorize an administrator and load the comment trash.
\newline -- \code{roleName}: \code{string} -- must equal \code{RoleName.ADMIN}.
\newline -- \code{page}, \code{limit}: \code{number} -- pagination.
\newline -- \textit{Returns:} \code{Promise<ForumTrashCommentListResponseType>} -- deleted comments and pagination.
\newline -- Throws \code{ForumReportForbiddenException} for a non-admin or a repository failure. \\ \hline

14 & \code{restoreTrashPost(}\newline\code{roleName: string, postId: number)} &
A method to authorize an administrator and restore a post.
\newline -- \code{roleName}: \code{string} -- must equal \code{RoleName.ADMIN}.
\newline -- \code{postId}: \code{number} -- post identifier.
\newline -- \textit{Returns:} \code{Promise<ForumPostType>} -- the restored post.
\newline -- Throws \code{FailedToRestoreForumPostException} for an unexpected failure. \\ \hline

15 & \code{restoreTrashComment(}\newline\code{roleName: string, commentId: number)} &
A method to authorize an administrator and restore a comment.
\newline -- \code{roleName}: \code{string} -- must equal \code{RoleName.ADMIN}.
\newline -- \code{commentId}: \code{number} -- comment identifier.
\newline -- \textit{Returns:} \code{Promise<ForumAdminCommentType>} -- the restored comment.
\newline -- A Prisma \code{P2025} error is mapped to \code{ForumCommentNotFoundException}; other Prisma errors become \code{FailedToRestoreForumCommentException}. \\ \hline
\end{ClassMethodTable}

\paragraph{ClassAttribute}
\begin{ClassAttributeTable}
01 & \code{adminRepository} & \code{ForumAdminRepository} -- private readonly persistence dependency. \\ \hline
\end{ClassAttributeTable}

\subsection{ForumAdminRepository Class}

\paragraph{ClassMethod}
\begin{ClassMethodTable}
01 & \code{ForumAdminRepository(prisma: PrismaService)} &
A method to initialize the repository with the shared Prisma service.
\newline -- \code{prisma}: \code{PrismaService} -- database client.
\newline -- \textit{Returns:} \code{void} -- constructor, no return value. \\ \hline

02 & \code{getPendingPosts(page: number, limit: number)} &
A method to query the moderation queue and its total in one transaction.
\newline -- \code{page}, \code{limit}: \code{number} -- pagination.
\newline -- \textit{Returns:} \code{Promise<{ posts: PendingForumPostType[]; total: number }>} -- non-deleted \code{PENDING} posts, oldest first, with author summary.
\newline -- \code{moderationCategories} is normalized to \code{string[] | null}. \\ \hline

03 & \code{reviewForumPost(}\newline\code{postId: number,}\newline\code{status: 'APPROVED' | 'REJECTED',}\newline\code{adminId: number)} &
A method to apply a moderation decision only to a non-deleted \code{PENDING} post.
\newline -- \code{postId}: \code{number} -- post identifier.
\newline -- \code{status}: \code{'APPROVED' | 'REJECTED'} -- target moderation status.
\newline -- \code{adminId}: \code{number} -- reviewer recorded in \code{reviewedById}.
\newline -- \textit{Returns:} \code{Promise<ForumPostType>} -- the reviewed post.
\newline -- \code{REJECTED} also sets \code{deletedAt}; a zero-row update throws \code{ForumPostNotFoundException}. \\ \hline

04 & \code{getAdminStats()} &
A method to count forum entities and load the five newest posts in one transaction.
\newline -- \textit{Returns:} \code{Promise<ForumAdminStatsType>} -- active category, post, comment, and user counts; total and pending reports; five recent posts.
\newline -- Posts, comments, categories, and users exclude soft-deleted rows; report totals include all rows. \\ \hline

05 & \code{getAdminCommentLists(}\newline\code{page: number, limit: number,}\newline\code{search?: string)} &
A method to query a page of active comments and the matching total in one transaction.
\newline -- \code{page}, \code{limit}: \code{number} -- pagination.
\newline -- \code{search}: \code{string | undefined} -- optional case-insensitive content filter.
\newline -- \textit{Returns:} \code{Promise<{ comments: ForumAdminCommentType[]; total: number }>} -- comments with author and parent-post summaries, newest first. \\ \hline

06 & \code{getAdminCategoryLists(}\newline\code{page: number, limit: number,}\newline\code{search?: string,}\newline\code{sortBy?: 'id' | 'name' | 'createdAt',}\newline\code{sortOrder?: 'asc' | 'desc',}\newline\code{deleted?: string)} &
A method to query a page of categories and the matching total in one transaction.
\newline -- \code{page}, \code{limit}: \code{number} -- pagination.
\newline -- \code{search}: \code{string | undefined} -- optional case-insensitive name filter.
\newline -- \code{sortBy}, \code{sortOrder}: optional sort controls; defaults are \code{'id'} and \code{'desc'}.
\newline -- \code{deleted}: \code{string | undefined} -- \code{'true'} selects only deleted rows, \code{'false'} only active rows, omitted selects both.
\newline -- \textit{Returns:} \code{Promise<{ categories: ForumAdminCategoryType[]; total: number }>} -- categories with \code{postCount} and pagination total. \\ \hline

07 & \code{getAdminCategoryById(id: number)} &
A method to load one active category with its public post count.
\newline -- \code{id}: \code{number} -- category identifier.
\newline -- \textit{Returns:} \code{Promise<ForumCategoryType | null>} -- the category detail or \code{null}. \\ \hline

08 & \code{createAdminCategory(data: { name: string; description?: string | null })} &
A method to create a category with a unique name and generated slug.
\newline -- \code{data.name}: \code{string} -- category name.
\newline -- \code{data.description}: \code{string | null | undefined} -- optional description.
\newline -- \textit{Returns:} \code{Promise<ForumCategoryType>} -- the created category with \code{postCount = 0}.
\newline -- Throws \code{ForumCategoryAlreadyExistsException} for a case-insensitive active-name conflict; slug collisions receive a numeric suffix. \\ \hline

09 & \code{updateAdminCategory(}\newline\code{id: number, data: UpdateForumCategoryBodyType)} &
A method to update a category and regenerate its slug when the name changes.
\newline -- \code{id}: \code{number} -- category identifier.
\newline -- \code{data}: \code{UpdateForumCategoryBodyType} -- optional name and description.
\newline -- \textit{Returns:} \code{Promise<ForumCategoryType>} -- the updated category.
\newline -- Throws \code{ForumCategoryNotFoundException} or \code{ForumCategoryAlreadyExistsException}. \\ \hline

10 & \code{deleteAdminCategory(id: number)} &
A method to detach posts and soft-delete a category in one transaction.
\newline -- \code{id}: \code{number} -- category identifier.
\newline -- \textit{Returns:} \code{Promise<{ message: string }>} -- deletion confirmation.
\newline -- Non-deleted posts in the category are reassigned to \code{categoryId = null} before the category receives \code{deletedAt}. \\ \hline

11 & \code{restoreAdminCategory(id: number)} &
A method to restore a soft-deleted category after a name-conflict check.
\newline -- \code{id}: \code{number} -- category identifier.
\newline -- \textit{Returns:} \code{Promise<ForumCategoryType>} -- the restored category.
\newline -- Throws \code{ForumCategoryNotFoundException} or \code{ForumCategoryAlreadyExistsException}. \\ \hline

12 & \code{getTrashPosts(page: number, limit: number)} &
A method to query deleted or rejected posts and the matching total in one transaction.
\newline -- \code{page}, \code{limit}: \code{number} -- pagination.
\newline -- \textit{Returns:} \code{Promise<{ posts: ForumTrashPostType[]; total: number }>} -- trash posts with \code{deletedAt} and author summary.
\newline -- Ordering is by \code{deletedAt} descending with null values last. \\ \hline

13 & \code{getTrashComments(page: number, limit: number)} &
A method to query soft-deleted comments and the matching total in one transaction.
\newline -- \code{page}, \code{limit}: \code{number} -- pagination.
\newline -- \textit{Returns:} \code{Promise<{ comments: ForumTrashCommentType[]; total: number }>} -- trash comments with \code{deletedAt} and parent-post summary, newest deletion first. \\ \hline

14 & \code{restoreTrashPost(postId: number)} &
A method to restore a post from the trash.
\newline -- \code{postId}: \code{number} -- post identifier.
\newline -- \textit{Returns:} \code{Promise<ForumPostType>} -- the restored post.
\newline -- A deleted \code{APPROVED} post becomes visible again; a \code{REJECTED} post returns to \code{PENDING} and its reviewer fields are cleared. \\ \hline

15 & \code{restoreTrashComment(commentId: number)} &
A method to clear \code{deletedAt} on a trashed comment.
\newline -- \code{commentId}: \code{number} -- comment identifier.
\newline -- \textit{Returns:} \code{Promise<ForumAdminCommentType>} -- the restored comment with author and parent-post summaries.
\newline -- Throws \code{ForumCommentNotFoundException} when no deleted row matches. \\ \hline
\end{ClassMethodTable}

\paragraph{ClassAttribute}
\begin{ClassAttributeTable}
01 & \code{prisma} & \code{PrismaService} -- private readonly database client. \\ \hline
\end{ClassAttributeTable}

\subsection{IdentityVerificationsAdminController Class}

\paragraph{ClassMethod}
\begin{ClassMethodTable}
01 & \code{IdentityVerificationsAdminController(}\newline\code{service: IdentityVerificationsAdminService)} &
A method to initialize the controller with the administration service.
\newline -- \code{service}: \code{IdentityVerificationsAdminService} -- service dependency.
\newline -- \textit{Returns:} \code{void} -- constructor, no return value. \\ \hline

02 & \code{list(query: IdentityVerificationAdminFilterDto)} &
A method to return a filtered, paginated document list.
\newline -- \code{query}: \code{IdentityVerificationAdminFilterDto} -- page, limit, optional status, and optional search.
\newline -- \textit{Returns:} \code{Promise<IdentityVerificationAdminListResponseType>} -- documents and pagination.
\newline -- Route: \code{GET /admin/identity-verifications}. \\ \hline

03 & \code{detail(id: number)} &
A method to return one document with user and reviewer detail.
\newline -- \code{id}: \code{number} -- document identifier parsed by \code{ParseIntPipe}.
\newline -- \textit{Returns:} \code{Promise<IdentityVerificationAdminDetailType>} -- document, user profile, freelancer verification state, and reviewing administrator.
\newline -- Route: \code{GET /admin/identity-verifications/:id}. \\ \hline

04 & \code{file(}\newline\code{id: number, side: string | undefined,}\newline\code{response: Response)} &
A method to stream a legacy \code{local://} document image.
\newline -- \code{id}: \code{number} -- document identifier.
\newline -- \code{side}: \code{string | undefined} -- \code{'back'} selects the back image; any other value selects the front.
\newline -- \code{response}: \code{Response} -- Express response used to set headers.
\newline -- \textit{Returns:} \code{Promise<StreamableFile>} -- file stream; sets inline \code{Content-Disposition} and \code{application/octet-stream}.
\newline -- Route: \code{GET /admin/identity-verifications/:id/file}. \\ \hline

05 & \code{approve(}\newline\code{id: number, adminId: number,}\newline\code{body: ApproveIdentityVerificationBodyDto)} &
A method to approve a pending document.
\newline -- \code{id}: \code{number} -- document identifier.
\newline -- \code{adminId}: \code{number} -- reviewing administrator from \code{@UserActive('userId')}.
\newline -- \code{body}: \code{ApproveIdentityVerificationBodyDto} -- optional \code{reviewNotes} up to 1000 characters.
\newline -- \textit{Returns:} \code{Promise<IdentityVerificationAdminActionResponseType>} -- the reviewed document detail.
\newline -- Route: \code{PATCH /admin/identity-verifications/:id/approve}. \\ \hline

06 & \code{reject(}\newline\code{id: number, adminId: number,}\newline\code{body: RejectIdentityVerificationBodyDto)} &
A method to reject a pending document.
\newline -- \code{id}: \code{number} -- document identifier.
\newline -- \code{adminId}: \code{number} -- reviewing administrator.
\newline -- \code{body}: \code{RejectIdentityVerificationBodyDto} -- optional trimmed \code{reviewNotes} up to 1000 characters.
\newline -- \textit{Returns:} \code{Promise<IdentityVerificationAdminActionResponseType>} -- the reviewed document detail.
\newline -- Route: \code{PATCH /admin/identity-verifications/:id/reject}. \\ \hline
\end{ClassMethodTable}

\paragraph{ClassAttribute}
\begin{ClassAttributeTable}
01 & \code{service} & \code{IdentityVerificationsAdminService} -- private readonly service dependency. \\ \hline
\end{ClassAttributeTable}

\subsection{IdentityVerificationsAdminService Class}

\paragraph{ClassMethod}
\begin{ClassMethodTable}
01 & \code{IdentityVerificationsAdminService(}\newline\code{repository: IdentityVerificationsAdminRepository)} &
A method to initialize the service with the identity-verification repository.
\newline -- \code{repository}: \code{IdentityVerificationsAdminRepository} -- persistence dependency.
\newline -- \textit{Returns:} \code{void} -- constructor, no return value. \\ \hline

02 & \code{presentFileUrl(document: T)} &
A private generic method to replace legacy \code{local://} file paths with API URLs.
\newline -- \code{document}: \code{T} -- object with \code{id}, \code{fileUrl}, and optional \code{backFileUrl}.
\newline -- \textit{Returns:} the same object with \code{fileUrl} set to the file route and, when needed, \code{backFileUrl} set to the same route with \code{?side=back}.
\newline -- Non-local URLs are returned unchanged. \\ \hline

03 & \code{presentList(documents: T[])} &
A private method to map \code{presentFileUrl} over a document list.
\newline -- \code{documents}: \code{T[]} -- list of document objects.
\newline -- \textit{Returns:} the mapped list. \\ \hline

04 & \code{list(params: { page: number; limit: number; status?: string; search?: string })} &
A method to load a filtered page of documents.
\newline -- \code{params}: object -- page, limit, optional status string, and optional search.
\newline -- \textit{Returns:} \code{Promise<IdentityVerificationAdminListResponseType>} -- documents with presentation URLs and pagination.
\newline -- \code{status} is cast to the Prisma \code{VerificationStatus} union before the repository call. \\ \hline

05 & \code{detail(id: number)} &
A method to load one document and present its file URLs.
\newline -- \code{id}: \code{number} -- document identifier.
\newline -- \textit{Returns:} \code{Promise<IdentityVerificationAdminDetailType>} -- the document detail.
\newline -- Throws \code{IdentityVerificationNotFoundException}. \\ \hline

06 & \code{review(}\newline\code{id: number, adminId: number,}\newline\code{status: ReviewableStatus,}\newline\code{reviewNotes: string | null)} &
A method to approve or reject a pending document.
\newline -- \code{id}: \code{number} -- document identifier.
\newline -- \code{adminId}: \code{number} -- reviewing administrator.
\newline -- \code{status}: \code{ReviewableStatus} -- \code{APPROVED} or \code{REJECTED}.
\newline -- \code{reviewNotes}: \code{string | null} -- optional reviewer note.
\newline -- \textit{Returns:} \code{Promise<IdentityVerificationAdminActionResponseType>} -- the reviewed document detail.
\newline -- Throws \code{IdentityVerificationNotFoundException} or \code{IdentityVerificationAlreadyReviewedException}. \\ \hline

07 & \code{getLocalFile(}\newline\code{id: number, side: 'front' | 'back' = 'front')} &
A method to resolve and validate a legacy local document file.
\newline -- \code{id}: \code{number} -- document identifier.
\newline -- \code{side}: \code{'front' | 'back'} -- which stored URL to use; default front.
\newline -- \textit{Returns:} \code{Promise<{ absolutePath: string; fileName: string }>} -- resolved path and base filename.
\newline -- Throws \code{IdentityVerificationNotFoundException} or \code{IdentityVerificationFileInvalidException}; rejects paths escaping the uploads root. \\ \hline
\end{ClassMethodTable}

\paragraph{ClassAttribute}
\begin{ClassAttributeTable}
01 & \code{repository} & \code{IdentityVerificationsAdminRepository} -- private readonly persistence dependency. \\ \hline
\end{ClassAttributeTable}

\subsection{IdentityVerificationsAdminRepository Class}

\paragraph{ClassMethod}
\begin{ClassMethodTable}
01 & \code{IdentityVerificationsAdminRepository(}\newline\code{prisma: PrismaService)} &
A method to initialize the repository with the shared Prisma service.
\newline -- \code{prisma}: \code{PrismaService} -- database client.
\newline -- \textit{Returns:} \code{void} -- constructor, no return value. \\ \hline

02 & \code{findDocuments(params: IdentityVerificationAdminListParams)} &
A method to query a filtered page of active documents and the matching total.
\newline -- \code{params}: \code{IdentityVerificationAdminListParams} -- page, limit, optional status, and optional search.
\newline -- \textit{Returns:} \code{Promise<{ documents: IdentityVerificationAdminListItemType[]; total: number }>} -- documents with user summary, newest first.
\newline -- Search matches user email or profile display name case-insensitively; \code{deletedAt} must be \code{null}. \\ \hline

03 & \code{findDocumentById(id: number)} &
A method to load the full detail projection of an active document.
\newline -- \code{id}: \code{number} -- document identifier.
\newline -- \textit{Returns:} \code{Promise<IdentityVerificationAdminDetailType | null>} -- document, reviewer, user, and freelancer summary, or \code{null}. \\ \hline

04 & \code{findDocumentRaw(id: number)} &
A method to load the minimal fields needed for review and file streaming.
\newline -- \code{id}: \code{number} -- document identifier.
\newline -- \textit{Returns:} \code{Promise<{ id: number; userId: number; fileUrl: string; backFileUrl: string | null; status: VerificationStatus } | null>} -- raw document fields, or \code{null}. \\ \hline

05 & \code{reviewDocument(}\newline\code{id: number, adminId: number,}\newline\code{status: VerificationStatus,}\newline\code{reviewNotes: string | null)} &
A method to persist a review and synchronize the freelancer verification flag in one transaction.
\newline -- \code{id}: \code{number} -- document identifier.
\newline -- \code{adminId}: \code{number} -- reviewing administrator.
\newline -- \code{status}: \code{VerificationStatus} -- \code{APPROVED} or \code{REJECTED}.
\newline -- \code{reviewNotes}: \code{string | null} -- reviewer note.
\newline -- \textit{Returns:} \code{Promise<IdentityVerificationAdminDetailType | null>} -- the updated detail, or \code{null}.
\newline -- Updates the document status, reviewer, notes, and \code{reviewedAt}; sets \code{FreelancerProfile.idVerified} to true only for \code{APPROVED}. Transaction limits are 10s wait and 15s timeout. \\ \hline
\end{ClassMethodTable}

\paragraph{ClassAttribute}
\begin{ClassAttributeTable}
01 & \code{prisma} & \code{PrismaService} -- private readonly database client. \\ \hline
\end{ClassAttributeTable}

\section{Database Design}
\subsection{Profile, Account and Review Persistence -- Nguyen Lu Nhat Huy}
The following model fields and constraints support the five sprint modules; names are taken from the Prisma schema. Related global database diagrams remain in the preceding sections.
\subsubsection{Profile}
\begin{longtable}{|L{3.3cm}|L{10.8cm}|}\hline
Field / constraint & Type and mapping \\ \hline
id & Int @id @default(autoincrement()) \\ \hline
userId & Int @unique \\ \hline
displayName & String? @db.VarChar(255) \\ \hline
avatarUrl & String? @db.VarChar(500) \\ \hline
coverUrl & String? @db.VarChar(500) \\ \hline
bio & String? \\ \hline
onlineStatus & Boolean @default(false) \\ \hline
availabilityStatus & AvailabilityStatus @default(OFFLINE) \\ \hline
profileCompletionPercent & Int @default(0) \\ \hline
walletAddress & String? @db.VarChar(42) \\ \hline
createdAt & DateTime @default(now()) \\ \hline
updatedAt & DateTime @updatedAt \\ \hline
user & User @relation(fields: [userId], references: [id], onDelete: Restrict, onUpdate: NoAction) \\ \hline
freelancerProfile & FreelancerProfile? \\ \hline
clientProfile & ClientProfile? \\ \hline
expertProfile & Expert? \\ \hline
socialLinks & SocialLink[] \\ \hline
revisions & ProfileRevision[] \\ \hline
\end{longtable}
\subsubsection{FreelancerProfile}
\begin{longtable}{|L{3.3cm}|L{10.8cm}|}\hline
Field / constraint & Type and mapping \\ \hline
id & Int @id @default(autoincrement()) \\ \hline
profileId & Int @unique \\ \hline
title & String? @db.VarChar(255) \\ \hline
education & String[] @default([]) \\ \hline
certifications & String[] @default([]) \\ \hline
languages & String[] @default([]) \\ \hline
cvUrl & String? @db.VarChar(500) \\ \hline
cvFileName & String? @db.VarChar(255) \\ \hline
cvPublicId & String? @db.VarChar(255) \\ \hline
cvUploadedAt & DateTime? \\ \hline
experience & String[] @default([]) \\ \hline
idVerified & Boolean @default(false) \\ \hline
createdAt & DateTime @default(now()) \\ \hline
updatedAt & DateTime @updatedAt \\ \hline
profile & Profile @relation(fields: [profileId], references: [id], onDelete: Restrict, onUpdate: NoAction) \\ \hline
skills & FreelancerSkill[] \\ \hline
portfolioItems & PortfolioItem[] \\ \hline
\end{longtable}
\subsubsection{FreelancerSkill}
\begin{longtable}{|L{3.3cm}|L{10.8cm}|}\hline
Field / constraint & Type and mapping \\ \hline
id & Int @id @default(autoincrement()) \\ \hline
freelancerProfileId & Int \\ \hline
skillId & Int \\ \hline
proficiencyLevel & Int @db.SmallInt // Range 1-10. \\ \hline
freelancerProfile & FreelancerProfile @relation(fields: [freelancerProfileId], references: [id], onDelete: Restrict, onUpdate: NoAction) \\ \hline
skill & Skill @relation(fields: [skillId], references: [id], onDelete: Restrict, onUpdate: NoAction) \\ \hline
@@unique([freelancerProfileId, & skillId]) \\ \hline
@@index([freelancerProfileId]) &  \\ \hline
@@index([skillId]) &  \\ \hline
\end{longtable}
\subsubsection{PortfolioItem}
\begin{longtable}{|L{3.3cm}|L{10.8cm}|}\hline
Field / constraint & Type and mapping \\ \hline
id & Int @id @default(autoincrement()) \\ \hline
freelancerProfileId & Int \\ \hline
title & String @db.VarChar(255) \\ \hline
description & String? \\ \hline
technologies & String[] @default([]) \\ \hline
projectUrl & String? @db.VarChar(500) \\ \hline
createdAt & DateTime @default(now()) \\ \hline
updatedAt & DateTime @updatedAt \\ \hline
deletedAt & DateTime? \\ \hline
freelancerProfile & FreelancerProfile @relation(fields: [freelancerProfileId], references: [id], onDelete: Restrict, onUpdate: NoAction) \\ \hline
likes & PortfolioLike[] \\ \hline
\end{longtable}
\subsubsection{ClientProfile}
\begin{longtable}{|L{3.3cm}|L{10.8cm}|}\hline
Field / constraint & Type and mapping \\ \hline
id & Int @id @default(autoincrement()) \\ \hline
profileId & Int @unique \\ \hline
companyName & String? @db.VarChar(255) \\ \hline
companyDescription & String? \\ \hline
website & String? @db.VarChar(255) \\ \hline
createdAt & DateTime @default(now()) \\ \hline
updatedAt & DateTime @updatedAt \\ \hline
profile & Profile @relation(fields: [profileId], references: [id], onDelete: Restrict, onUpdate: NoAction) \\ \hline
\end{longtable}
\subsubsection{Expert}
\begin{longtable}{|L{3.3cm}|L{10.8cm}|}\hline
Field / constraint & Type and mapping \\ \hline
id & Int @id @default(autoincrement()) \\ \hline
profileId & Int @unique \\ \hline
title & String? @db.VarChar(255) \\ \hline
expertise & String[] @default([]) \\ \hline
yearsOfExperience & Int @default(0) \\ \hline
education & String[] @default([]) \\ \hline
certifications & String[] @default([]) \\ \hline
website & String? @db.VarChar(500) \\ \hline
isActive & Boolean @default(false) \\ \hline
createdAt & DateTime @default(now()) \\ \hline
updatedAt & DateTime @updatedAt \\ \hline
profile & Profile @relation(fields: [profileId], references: [id], onDelete: Restrict, onUpdate: NoAction) \\ \hline
@@index([isActive]) &  \\ \hline
\end{longtable}
\subsubsection{SocialLink}
\begin{longtable}{|L{3.3cm}|L{10.8cm}|}\hline
Field / constraint & Type and mapping \\ \hline
id & Int @id @default(autoincrement()) \\ \hline
profileId & Int \\ \hline
platform & SocialPlatform \\ \hline
url & String @db.VarChar(500) \\ \hline
profile & Profile @relation(fields: [profileId], references: [id], onDelete: Restrict, onUpdate: NoAction) \\ \hline
\end{longtable}
\subsubsection{FavoriteFreelancer}
\begin{longtable}{|L{3.3cm}|L{10.8cm}|}\hline
Field / constraint & Type and mapping \\ \hline
clientId & Int \\ \hline
freelancerId & Int \\ \hline
createdAt & DateTime @default(now()) \\ \hline
client & User @relation("FavoriteClient", fields: [clientId], references: [id], onDelete: Cascade, onUpdate: NoAction) \\ \hline
freelancer & User @relation("FavoriteTargetFreelancer", fields: [freelancerId], references: [id], onDelete: Cascade, onUpdate: NoAction) \\ \hline
@@id([clientId, & freelancerId]) \\ \hline
@@index([freelancerId]) &  \\ \hline
\end{longtable}
\subsubsection{FollowFreelancer}
\begin{longtable}{|L{3.3cm}|L{10.8cm}|}\hline
Field / constraint & Type and mapping \\ \hline
clientId & Int \\ \hline
freelancerId & Int \\ \hline
createdAt & DateTime @default(now()) \\ \hline
client & User @relation("FollowClient", fields: [clientId], references: [id], onDelete: Cascade, onUpdate: NoAction) \\ \hline
freelancer & User @relation("FollowTargetFreelancer", fields: [freelancerId], references: [id], onDelete: Cascade, onUpdate: NoAction) \\ \hline
@@id([clientId, & freelancerId]) \\ \hline
@@index([freelancerId]) &  \\ \hline
\end{longtable}
\subsubsection{Notification}
\begin{longtable}{|L{3.3cm}|L{10.8cm}|}\hline
Field / constraint & Type and mapping \\ \hline
id & Int @id @default(autoincrement()) \\ \hline
userId & Int \\ \hline
type & NotificationType \\ \hline
title & String? @db.VarChar(255) \\ \hline
message & String? \\ \hline
data & Json? \\ \hline
isRead & Boolean @default(false) \\ \hline
createdAt & DateTime @default(now()) \\ \hline
deletedAt & DateTime? \\ \hline
user & User @relation(fields: [userId], references: [id], onDelete: Restrict, onUpdate: NoAction) \\ \hline
@@index([userId, & isRead]) \\ \hline
@@index([createdAt]) &  \\ \hline
\end{longtable}
\subsubsection{Review}
\begin{longtable}{|L{3.3cm}|L{10.8cm}|}\hline
Field / constraint & Type and mapping \\ \hline
id & Int @id @default(autoincrement()) \\ \hline
contractId & Int \\ \hline
reviewerId & Int \\ \hline
revieweeId & Int \\ \hline
overallRating & Decimal @db.Decimal(2, 1) \\ \hline
breakdown & Json? \\ \hline
comment & String? \\ \hline
createdAt & DateTime @default(now()) \\ \hline
updatedAt & DateTime @updatedAt \\ \hline
deletedAt & DateTime? \\ \hline
contract & Contract @relation(fields: [contractId], references: [id], onDelete: Restrict, onUpdate: NoAction) \\ \hline
reviewer & User @relation("ReviewReviewer", fields: [reviewerId], references: [id], onDelete: Restrict, onUpdate: NoAction) \\ \hline
reviewee & User @relation("ReviewReviewee", fields: [revieweeId], references: [id], onDelete: Restrict, onUpdate: NoAction) \\ \hline
response & ReviewResponse? \\ \hline
@@unique([contractId, & reviewerId]) \\ \hline
@@index([revieweeId]) &  \\ \hline
@@index([reviewerId]) &  \\ \hline
\end{longtable}
\subsubsection{ReviewResponse}
\begin{longtable}{|L{3.3cm}|L{10.8cm}|}\hline
Field / constraint & Type and mapping \\ \hline
id & Int @id @default(autoincrement()) \\ \hline
reviewId & Int @unique \\ \hline
userId & Int \\ \hline
responseText & String \\ \hline
createdAt & DateTime @default(now()) \\ \hline
updatedAt & DateTime @updatedAt \\ \hline
deletedAt & DateTime? \\ \hline
review & Review @relation(fields: [reviewId], references: [id], onDelete: Restrict, onUpdate: NoAction) \\ \hline
user & User @relation(fields: [userId], references: [id], onDelete: Restrict, onUpdate: NoAction) \\ \hline
\end{longtable}
\subsubsection{ProfileRevision}
\begin{longtable}{|L{3.3cm}|L{10.8cm}|}\hline
Field / constraint & Type and mapping \\ \hline
id & Int @id @default(autoincrement()) \\ \hline
userId & Int \\ \hline
profileId & Int \\ \hline
profileType & ProfileRevisionType \\ \hline
status & RevisionStatus @default(PENDING) \\ \hline
currentData & Json \\ \hline
proposedData & Json \\ \hline
profileStrength & Int \\ \hline
adminId & Int? \\ \hline
reviewNotes & String? \\ \hline
createdAt & DateTime @default(now()) \\ \hline
updatedAt & DateTime @updatedAt \\ \hline
reviewedAt & DateTime? \\ \hline
user & User @relation("ProfileRevisionUser", fields: [userId], references: [id], onDelete: Restrict, onUpdate: NoAction) \\ \hline
admin & User? @relation("ProfileRevisionAdmin", fields: [adminId], references: [id], onDelete: SetNull, onUpdate: NoAction) \\ \hline
profile & Profile @relation(fields: [profileId], references: [id], onDelete: Restrict, onUpdate: NoAction) \\ \hline
@@index([userId, & profileType, status]) \\ \hline
@@index([status, & createdAt]) \\ \hline
@@index([profileId]) &  \\ \hline
\end{longtable}
\subsubsection{IdVerificationDocument}
\begin{longtable}{|L{3.3cm}|L{10.8cm}|}\hline
Field / constraint & Type and mapping \\ \hline
id & Int @id @default(autoincrement()) \\ \hline
userId & Int \\ \hline
documentType & DocumentType \\ \hline
fileUrl & String @db.VarChar(500) \\ \hline
status & VerificationStatus @default(PENDING) \\ \hline
adminId & Int? \\ \hline
reviewNotes & String? \\ \hline
createdAt & DateTime @default(now()) \\ \hline
reviewedAt & DateTime? \\ \hline
deletedAt & DateTime? \\ \hline
aiStatus & String? @db.VarChar(20) \\ \hline
aiScore & Float? \\ \hline
aiSummary & Json? \\ \hline
aiSelfieSummary & Json? \\ \hline
aiError & String? \\ \hline
aiScannedAt & DateTime? \\ \hline
backFileUrl & String? @db.VarChar(500) \\ \hline
identityNumber & String? @db.VarChar(50) \\ \hline
duplicateOfId & Int? \\ \hline
user & User @relation("IdVerificationUser", fields: [userId], references: [id], onDelete: Restrict, onUpdate: NoAction) \\ \hline
admin & User? @relation("IdVerificationAdmin", fields: [adminId], references: [id], onDelete: SetNull, onUpdate: NoAction) \\ \hline
duplicateOf & IdVerificationDocument? @relation("IdVerificationDuplicate", fields: [duplicateOfId], references: [id], onDelete: SetNull, onUpdate: NoAction) \\ \hline
duplicates & IdVerificationDocument[] @relation("IdVerificationDuplicate") \\ \hline
@@index([identityNumber]) &  \\ \hline
\end{longtable}
\subsubsection{UserRole}
\begin{longtable}{|L{3.3cm}|L{10.8cm}|}\hline
Field / constraint & Type and mapping \\ \hline
userId & Int \\ \hline
roleId & Int \\ \hline
isPrimary & Boolean @default(false) \\ \hline
createdAt & DateTime @default(now()) \\ \hline
user & User @relation(fields: [userId], references: [id], onDelete: Cascade, onUpdate: NoAction) \\ \hline
role & Role @relation(fields: [roleId], references: [id], onDelete: Cascade, onUpdate: NoAction) \\ \hline
@@id([userId, & roleId]) \\ \hline
\end{longtable}



\end{document}
